% supplement.tex: electronic supplement to main.tex (Paper 1, ACM Computing Surveys).
% Compiled separately: latexmk -pdf supplement.tex (main.tex is compiled first so that cross-references to it resolve; xr-hyper reads main.aux).
% Not counted in the 35 pages of the main paper (CSUR author guidelines, "Paper Length"; research/review/dossier/csur.md Part A).
% Sections keep the numbers promised in the main text and in OUTLINE.md section 1 (S1 search strings, S3 census tables, S5 design-space table in full,
% S7 EXP-003 detail). Material moved out of the main text on 2026-10-15 to bring it to 35 pages is in S5, S7 and S8; every moved block keeps its
% original "% src:" comment. S1 to S4 and S6 of the OUTLINE are not assembled yet.
\documentclass[acmsmall]{acmart}

\AtBeginDocument{\providecommand\BibTeX{{Bib\TeX}}}
\citestyle{acmauthoryear}

\usepackage[ruled,vlined,linesnumbered]{algorithm2e}
\usepackage{longtable}
% Cross-references to main.tex: scripts/mainrefs.py reads main.aux and writes mainrefs.tex (run it after compiling main.tex).
\makeatletter
\InputIfFileExists{mainrefs.tex}{}{}
\newcommand{\Mref}[1]{\@ifundefined{mr@#1}{\textbf{??}}{\csname mr@#1\endcsname}}
\makeatother
\newcommand{\mask}{\mathsf{m}}                 % mask (erasure) symbol, not in the vocabulary
\newcommand{\KL}[2]{\mathrm{KL}\!\left(#1\,\middle\|\,#2\right)}
\newcommand{\TC}{\mathrm{TC}}                  % total correlation
\newcommand{\Efact}{E_{\mathrm{fact}}}
\newcommand{\Elearn}{E_{\mathrm{learn}}}
\newcommand{\Eapprox}{E_{\mathrm{approx}}}

\acmJournal{CSUR}
\acmVolume{1}\acmNumber{1}\acmArticle{1}\acmMonth{1}\acmYear{2027}

% ---- drafting aids (same as main.tex) ----
\newcommand{\todo}[1]{{\color{red}\textbf{[TODO: #1]}}}
\newcommand{\aud}{\textsuperscript{\color{red}\scriptsize audit}}

% Supplement numbering: sections S5, S7, ...; tables S1, ...; algorithms S1, ...
\renewcommand{\thesection}{S\arabic{section}}
\renewcommand{\thetable}{S\arabic{table}}
\renewcommand{\thealgocf}{S\arabic{algocf}}

\begin{document}

\title[Supplement to the review of diffusion language models]{Supplement to ``Diffusion Language Models and Unified Multimodal Generation: A Critical Systematic Review of Discrete and Continuous Approaches''}

\author{Shlok Kumar Goenka}
\affiliation{%
  \institution{Vellore Institute of Technology}
  \city{Vellore}
  \state{Tamil Nadu}
  \country{India}}
\email{shlok.goenka2023@vitstudent.ac.in}

\begin{abstract}
This file holds the material that the main paper points to and that is not counted in its 35 pages: the condensed design-space table of the 47 models, the details of the EXP-003 measurement that are not needed to follow the argument, and four passages moved out of Sections 3, 4 and 7 of the main text. Section numbers refer to this file when written S\textit{n}, and to the main paper otherwise.
\end{abstract}

\maketitle

\todo{assemble S2 (screening table), S3 (census codebooks and the 176-row table), S4 (agreement tables), S6 (extraction and quality-appraisal tables) and the readme.txt that CSUR asks for; OUTLINE.md section 1 lists them}

\setcounter{section}{0}
\section{Search strings and PRISMA 2020 checklist}
\label{sec:s-search}

\subsection{PRISMA 2020 checklist}
Item numbers and topics follow the PRISMA 2020 statement~\cite{page2021prisma}; ``Location'' gives the section of the main paper (Section~\textit{n}) or of this supplement (S\textit{n}). Where the review meets an item only in part, the note says so.
{\small
\begin{longtable}{@{}p{1.2cm}p{3.6cm}p{7.8cm}@{}}
\toprule
Item & Topic & Location and note \\
\midrule
\endhead
1 & Title & Title page: ``systematic review'' in the title. \\
2 & Abstract & Abstract (journal limit 100 words; the PRISMA abstract items are partly covered). \\
3 & Rationale & Section~\Mref{sec:intro-why}. \\
4 & Objectives & Section~\Mref{sec:intro-rq} (RQ1 to RQ6). \\
5 & Eligibility criteria & Section~\Mref{sec:method-eligibility}. \\
6 & Information sources & Section~\Mref{sec:method-sources} and Table~\Mref{tab:sources}, with dates. \\
7 & Search strategy & S1.2: every logged query, verbatim. \\
8 & Selection process & Section~\Mref{sec:method-screening}: two model screeners and a model adjudicator; human validation not done. \\
9 & Data collection process & Section~\Mref{sec:method-extraction}: one model extractor, a second on 8 studies, a validator. \\
10a, 10b & Data items & Section~\Mref{sec:method-extraction}; the extraction schema in the repository. \\
11 & Study risk of bias & Section~\Mref{sec:method-extraction}: QA1 to QA9, an appraisal of reporting and reproducibility, not a risk-of-bias tool for trials. \\
12 & Effect measures & Section~\Mref{sec:rq}: direction of within-paper contrasts; no pooled effect. \\
13a--13f & Synthesis methods & Section~\Mref{sec:method-extraction} (``Synthesis''): a verdict per question with stratum weights; no meta-analysis, so heterogeneity (13e) is not assessed. \\
14 & Reporting bias assessment & Not assessed formally; discussed in Section~\Mref{sec:limits}. \\
15 & Certainty assessment & Verdict vocabulary (Section~\Mref{sec:method-extraction}); no GRADE-type rating. \\
16a & Study selection & Figure~\Mref{fig:prisma}. \\
16b & Excluded studies & Exclusion code for every full-text exclusion (Figure~\Mref{fig:prisma}); per-record list in the repository. \\
17 & Study characteristics & Section~\Mref{sec:rq} and S6. \\
18 & Risk of bias in studies & Section~\Mref{sec:rq} and S6. \\
19 & Results of individual studies & S6. \\
20a--20d & Results of syntheses & Section~\Mref{sec:rq}. \\
21 & Reporting biases & Section~\Mref{sec:limits}. \\
22 & Certainty of evidence & Section~\Mref{sec:rq} (verdicts). \\
23a--23d & Discussion & Sections~\Mref{sec:open}, \Mref{sec:limits} and \Mref{sec:conclusion}. \\
24a & Registration & Not registered (Section~\Mref{sec:method-design}). \\
24b & Protocol & Repository file \texttt{research/review/protocol.md}. \\
24c & Amendments & Section~\Mref{sec:method-deviations}; decision ledger in the repository. \\
25, 26 & Support; competing interests & Acknowledgements. \\
27 & Availability of data and code & Repository and archive (Section~\Mref{sec:ai-use}). \\
\bottomrule
\end{longtable}
}

\subsection{Every logged search}
Generated from the search log by \texttt{paper/review/scripts/search\_strings\_table.py}; the exports and their SHA-256 values are in the repository.

% generated by paper/review/scripts/search_strings_table.py from research/review/corpus/search-log.jsonl; do not edit
\begin{small}
\paragraph{arXiv API (export.arxiv.org)}
\noindent\textbf{S1}, 2026-10-14 07:48 UTC, 1158 records.\\
{\ttfamily\footnotesize ((abs:"diffusion language model" OR abs:"diffusion language models" OR abs:"discrete diffusion" OR abs:"masked diffusion" OR abs:"absorbing diffusion" OR abs:"discrete flow matching" OR abs:"diffusion LLM" OR abs:"diffusion large language model" OR abs:"text diffusion" OR ti:"Diffusion-LM") AND (abs:language OR abs:text OR abs:token OR abs:tokens)) AND (cat:cs.CL OR cat:cs.LG OR cat:cs.CV OR cat:cs.AI OR cat:stat.ML) AND submittedDate:[202101010000 TO 202610142359]}\par\smallskip
\noindent\textbf{S2}, 2026-10-14 07:49 UTC, 173 records.\\
{\ttfamily\footnotesize ((abs:"unified multimodal" OR abs:"understanding and generation" OR abs:"any-to-any") AND (abs:diffusion OR abs:"masked generative" OR abs:"discrete flow" OR abs:"flow matching") AND (abs:multimodal OR abs:"vision-language" OR abs:"text-to-image")) AND (cat:cs.CL OR cat:cs.LG OR cat:cs.CV OR cat:cs.AI OR cat:stat.ML) AND submittedDate:[202101010000 TO 202610142359]}\par\smallskip
\paragraph{OpenAlex works API (title\_and\_abstract.search)}
\noindent\textbf{S2}, 2026-10-15 10:55 UTC, 115 records.\\
{\ttfamily\footnotesize title\_and\_abstract.search:("unified multimodal" OR "unified understanding and generation" OR "understanding and generation" OR "any-to-any") AND (diffusion OR "masked generative" OR "discrete flow" OR "flow matching") AND (multimodal OR "vision-language" OR "text-to-image"),from\_publication\_date:2021-01-01,to\_publication\_date:2026-10-15,language:en}\par\smallskip
\noindent\textbf{S1}, 2026-10-15 10:55 UTC, 200 records (hits reported 566).\\
{\ttfamily\footnotesize title\_and\_abstract.search:("diffusion language model" OR "discrete diffusion" OR "masked diffusion" OR "absorbing diffusion" OR "absorbing state" OR "discrete flow matching" OR "diffusion LLM" OR "diffusion large language model" OR "text diffusion" OR ("non-autoregressive" AND diffusion)) AND (language OR text OR token OR "language model"),from\_publication\_date:2021-01-01,to\_publication\_date:2026-10-15,language:en}\par\smallskip
\noindent\textbf{S1}, 2026-10-15 10:56 UTC, 566 records.\\
{\ttfamily\footnotesize title\_and\_abstract.search:("diffusion language model" OR "discrete diffusion" OR "masked diffusion" OR "absorbing diffusion" OR "absorbing state" OR "discrete flow matching" OR "diffusion LLM" OR "diffusion large language model" OR "text diffusion" OR ("non-autoregressive" AND diffusion)) AND (language OR text OR token OR "language model"),from\_publication\_date:2021-01-01,to\_publication\_date:2026-10-15,language:en}\par\smallskip
\noindent\textbf{S1':A1}, 2026-10-15 12:55 UTC, 200 records (hits reported 412).\\
{\ttfamily\footnotesize title\_and\_abstract.search:"diffusion language model" OR "language diffusion model" OR "diffusion llm" OR "diffusion lm" OR dllm OR "diffusion-based language model" OR "diffusion-based llm" OR "diffusion large language model" OR "diffusion-based large language model" OR "text diffusion" OR "block diffusion" OR "absorbing diffusion" OR "diffusion mllm" OR "diffusion vlm" OR "diffusion vision-language model" OR dmllm OR dvlm,from\_publication\_date:2021-01-01,to\_publication\_date:2026-10-15,language:en}\par\smallskip
\noindent\textbf{S1':A2}, 2026-10-15 12:55 UTC, 200 records (hits reported 925).\\
{\ttfamily\footnotesize title\_and\_abstract.search:(("discrete diffusion" OR "masked diffusion" OR "discrete flow" OR "discrete state-space") AND (language OR text OR (token OR tokens OR tokenizer) OR "language model" OR llm OR vocabulary OR autoregressive OR reasoning)) OR (("discrete diffusion" OR "masked diffusion" OR "discrete flow" OR "discrete state-space") AND NOT (protein OR molecule OR molecular OR antibody OR peptide OR rna OR dna OR drug OR crystal OR chemical OR chemistry OR genome OR genomic OR biology OR biological OR "graph generation" OR "traveling salesman" OR recommendation OR recommender) AND NOT (image OR video OR 3d OR audio OR visual OR visually OR pixel OR photo OR photograph OR segmentation)),from\_publication\_date:2021-01-01,to\_publication\_date:2026-10-15,language:en}\par\smallskip
\noindent\textbf{S1':A3}, 2026-10-15 13:04 UTC, 0 records.\\
{\ttfamily\footnotesize title\_and\_abstract.search:("absorbing state" AND (language OR text OR (token OR tokens OR tokenizer) OR "language model")) OR ("any-order" AND (language OR text OR (token OR tokens OR tokenizer) OR "language model")) OR "non-autoregressive language model" OR ("non-autoregressive" AND diffusion AND NOT (protein OR molecule OR molecular OR antibody OR peptide OR rna OR dna OR drug OR crystal OR chemical OR chemistry OR genome OR genomic OR biology OR biological OR "graph generation" OR "traveling salesman" OR recommendation OR recommender)),from\_publication\_date:2021-01-01,to\_publication\_date:2026-10-15,language:en}\par\smallskip
\noindent\textbf{S1':B}, 2026-10-15 13:04 UTC, 0 records.\\
{\ttfamily\footnotesize title\_and\_abstract.search:diffusion AND ("text generation" OR "language generation" OR "natural language generation" OR "sequence-to-sequence" OR seq2seq OR "discrete sequence" OR "speculative decoding" OR "categorical data" OR "discrete data" OR (paraphrase OR paraphrasing) OR "textual diffusion" OR "text modeling" OR "named entity" OR sentiment OR summarization OR (dialog OR dialogue) OR dialogue OR "machine translation" OR "text style" OR commonsense OR "topic model" OR "triple extraction" OR "information extraction") AND NOT (protein OR molecule OR molecular OR antibody OR peptide OR rna OR dna OR drug OR crystal OR chemical OR chemistry OR genome OR genomic OR biology OR biological OR "graph generation" OR "traveling salesman" OR recommendation OR recommender),from\_publication\_date:2021-01-01,to\_publication\_date:2026-10-15,language:en}\par\smallskip
\noindent\textbf{S1':C}, 2026-10-15 13:04 UTC, 0 records.\\
{\ttfamily\footnotesize title\_and\_abstract.search:diffusion AND ("language model" OR llm) AND NOT (protein OR molecule OR molecular OR antibody OR peptide OR rna OR dna OR drug OR crystal OR chemical OR chemistry OR genome OR genomic OR biology OR biological OR "graph generation" OR "traveling salesman" OR recommendation OR recommender) AND NOT (image OR video OR 3d OR audio OR visual OR visually OR pixel OR photo OR photograph OR segmentation),from\_publication\_date:2021-01-01,to\_publication\_date:2026-10-15,language:en}\par\smallskip
\noindent\textbf{S2':M}, 2026-10-15 13:04 UTC, 0 records.\\
{\ttfamily\footnotesize title\_and\_abstract.search:("unified multimodal" OR "unified multi-modal" OR "unified understanding and generation" OR "understanding and generation" OR "any-to-any" OR "mixed-modal" OR "mixed-modality") AND (diffusion OR "masked generative" OR "discrete flow" OR "flow matching" OR "rectified flow" OR "discrete denoising") AND (multimodal OR "multi-modal" OR "vision-language" OR "text-to-image" OR "image generation" OR mllm OR "visual understanding" OR "image understanding"),from\_publication\_date:2021-01-01,to\_publication\_date:2026-10-15,language:en}\par\smallskip
\noindent\textbf{S2':VLA}, 2026-10-15 13:04 UTC, 0 records.\\
{\ttfamily\footnotesize title\_and\_abstract.search:("vision-language-action" OR vla) AND ("discrete diffusion" OR "discrete denoising" OR "masked diffusion" OR "diffusion-based" OR "unified diffusion" OR "diffusion language model" OR dllm),from\_publication\_date:2021-01-01,to\_publication\_date:2026-10-15,language:en}\par\smallskip
\paragraph{Semantic Scholar Graph API (paper/search/bulk)}
\noindent\textbf{S1}, 2026-10-15 10:55 UTC, 1487 records.\\
{\ttfamily\footnotesize ("diffusion language model*" | "discrete diffusion" | "masked diffusion" | "absorbing diffusion" | "absorbing state" | "discrete flow matching" | "diffusion LLM" | "diffusion large language model*" | "text diffusion" | ("non-autoregressive" + diffusion)) + (language | text | token* | "language model*")}\par\smallskip
\noindent\textbf{S2}, 2026-10-15 10:56 UTC, 228 records.\\
{\ttfamily\footnotesize ("unified multimodal" | "unified understanding and generation" | "understanding and generation" | "any-to-any") + (diffusion | "masked generative" | "discrete flow" | "flow matching") + (multimodal | "vision-language" | "text-to-image")}\par\smallskip
\noindent\textbf{S1':A1}, 2026-10-15 12:55 UTC, 1769 records.\\
{\ttfamily\footnotesize "diffusion language model*" | "language diffusion model*" | "diffusion llm*" | "diffusion lm*" | dllm* | "diffusion-based language model*" | "diffusion-based llm*" | "diffusion large language model*" | "diffusion-based large language model*" | "text diffusion" | "block diffusion" | "absorbing diffusion" | "diffusion mllm*" | "diffusion vlm*" | "diffusion vision-language model*" | dmllm* | dvlm*}\par\smallskip
\noindent\textbf{S1':A2}, 2026-10-15 13:11 UTC, 1346 records.\\
{\ttfamily\footnotesize (("discrete diffusion" | "masked diffusion" | "discrete flow*" | "discrete state-space*") + (language | text | token* | "language model*" | llm* | vocabulary | autoregressive | reasoning)) | (("discrete diffusion" | "masked diffusion" | "discrete flow*" | "discrete state-space*") + -protein* + -molecul* + -antibod* + -peptide* + -rna + -dna + -drug* + -crystal* + -chemi* + -genom* + -biolog* + -"graph generation" + -"traveling salesman" + -recommend* + -image* + -video* + -3d + -audio + -visual* + -pixel* + -photo* + -segmentation)}\par\smallskip
\noindent\textbf{S1':A3}, 2026-10-15 13:15 UTC, 354 records.\\
{\ttfamily\footnotesize ("absorbing state" + (language | text | token* | "language model*")) | ("any-order" + (language | text | token* | "language model*")) | "non-autoregressive language model*" | ("non-autoregressive" + diffusion + -protein* + -molecul* + -antibod* + -peptide* + -rna + -dna + -drug* + -crystal* + -chemi* + -genom* + -biolog* + -"graph generation" + -"traveling salesman" + -recommend*)}\par\smallskip
\noindent\textbf{S1':B}, 2026-10-15 13:16 UTC, 3677 records.\\
{\ttfamily\footnotesize diffusion + ("text generation" | "language generation" | "natural language generation" | "sequence-to-sequence" | seq2seq | "discrete sequence*" | "speculative decoding" | "categorical data" | "discrete data" | paraphras* | "textual diffusion" | "text modeling" | "named entity" | sentiment | summarization | dialog* | dialogue | "machine translation" | "text style" | commonsense | "topic model*" | "triple extraction" | "information extraction") + -protein* + -molecul* + -antibod* + -peptide* + -rna + -dna + -drug* + -crystal* + -chemi* + -genom* + -biolog* + -"graph generation" + -"traveling salesman" + -recommend*}\par\smallskip
\noindent\textbf{S1':C}, 2026-10-15 13:24 UTC, 1914 records.\\
{\ttfamily\footnotesize diffusion + ("language model*" | llm*) + -protein* + -molecul* + -antibod* + -peptide* + -rna + -dna + -drug* + -crystal* + -chemi* + -genom* + -biolog* + -"graph generation" + -"traveling salesman" + -recommend* + -image* + -video* + -3d + -audio + -visual* + -pixel* + -photo* + -segmentation}\par\smallskip
\noindent\textbf{S2':M}, 2026-10-15 13:25 UTC, 289 records.\\
{\ttfamily\footnotesize ("unified multimodal" | "unified multi-modal" | "unified understanding and generation" | "understanding and generation" | "any-to-any" | "mixed-modal*" | "mixed-modality") + (diffusion | "masked generative" | "discrete flow*" | "flow matching" | "rectified flow" | "discrete denoising") + (multimodal | "multi-modal" | "vision-language" | "text-to-image" | "image generation" | mllm* | "visual understanding" | "image understanding")}\par\smallskip
\noindent\textbf{S2':VLA}, 2026-10-15 13:29 UTC, 111 records.\\
{\ttfamily\footnotesize ("vision-language-action" | vla*) + ("discrete diffusion" | "discrete denoising" | "masked diffusion" | "diffusion-based" | "unified diffusion" | "diffusion language model*" | dllm*)}\par\smallskip
\noindent\textbf{S1':A1}, 2026-10-15 13:36 UTC, 1171 records.\\
{\ttfamily\footnotesize "diffusion language model*" | "language diffusion model*" | "diffusion llm*" | "diffusion lm*" | dllm | dllms | "diffusion-based language model*" | "diffusion-based llm*" | "diffusion large language model*" | "diffusion-based large language model*" | "text diffusion" | "block diffusion" | "absorbing diffusion" | "diffusion mllm*" | "diffusion vlm*" | "diffusion vision-language model*" | dmllm | dmllms | dvlm | dvlms}\par\smallskip
\noindent\textbf{S2':M}, 2026-10-15 13:36 UTC, 289 records.\\
{\ttfamily\footnotesize ("unified multimodal" | "unified multi-modal" | "unified understanding and generation" | "understanding and generation" | "any-to-any" | "mixed-modal*" | "mixed-modality") + (diffusion | "masked generative" | "discrete flow*" | "flow matching" | "rectified flow" | "discrete denoising") + (multimodal | "multi-modal" | "vision-language" | "text-to-image" | "image generation" | mllm | mllms | "visual understanding" | "image understanding")}\par\smallskip
\noindent\textbf{S2':VLA}, 2026-10-15 13:36 UTC, 106 records.\\
{\ttfamily\footnotesize ("vision-language-action" | vla | vlas) + ("discrete diffusion" | "discrete denoising" | "masked diffusion" | "diffusion-based" | "unified diffusion" | "diffusion language model*" | dllm | dllms)}\par\smallskip
\paragraph{arXiv API (export.arxiv.org), protocol string verbatim}
\noindent\textbf{S1}, 2026-10-15 10:55 UTC, 1286 records.\\
{\ttfamily\footnotesize (((ti:"diffusion language model" OR abs:"diffusion language model") OR (ti:"diffusion language models" OR abs:"diffusion language models") OR (ti:"discrete diffusion" OR abs:"discrete diffusion") OR (ti:"masked diffusion" OR abs:"masked diffusion") OR (ti:"absorbing diffusion" OR abs:"absorbing diffusion") OR (ti:"absorbing state" OR abs:"absorbing state") OR (ti:"discrete flow matching" OR abs:"discrete flow matching") OR (ti:"diffusion LLM" OR abs:"diffusion LLM") OR (ti:"diffusion large language model" OR abs:"diffusion large language model") OR (ti:"diffusion large language models" OR abs:"diffusion large language models") OR (ti:"text diffusion" OR abs:"text diffusion") OR ((ti:"non-autoregressive" OR abs:"non-autoregressive") AND (ti:diffusion OR abs:diffusion))) AND ((ti:language OR abs:language) OR (ti:text OR abs:text) OR (ti:token* OR abs:token*) OR (ti:"language model" OR abs:"language model") OR (ti:"language models" OR abs:"language models"))) AND (cat:cs.CL OR cat:cs.LG OR cat:cs.CV OR cat:cs.AI OR cat:stat.ML) AND submittedDate:[202101010000 TO 202610152359]}\par\smallskip
\noindent\textbf{S2}, 2026-10-15 10:56 UTC, 190 records.\\
{\ttfamily\footnotesize (((ti:"unified multimodal" OR abs:"unified multimodal") OR (ti:"unified understanding and generation" OR abs:"unified understanding and generation") OR (ti:"understanding and generation" OR abs:"understanding and generation") OR (ti:"any-to-any" OR abs:"any-to-any")) AND ((ti:diffusion OR abs:diffusion) OR (ti:"masked generative" OR abs:"masked generative") OR (ti:"discrete flow" OR abs:"discrete flow") OR (ti:"flow matching" OR abs:"flow matching")) AND ((ti:multimodal OR abs:multimodal) OR (ti:"vision-language" OR abs:"vision-language") OR (ti:"text-to-image" OR abs:"text-to-image"))) AND (cat:cs.CL OR cat:cs.LG OR cat:cs.CV OR cat:cs.AI OR cat:stat.ML) AND submittedDate:[202101010000 TO 202610152359]}\par\smallskip
\paragraph{ACL Anthology (metadata export anthology+abstracts.bib.gz, string applied locally)}
\noindent\textbf{S1}, 2026-10-15 10:55 UTC, 113 records.\\
{\ttfamily\footnotesize ("diffusion language model*" OR "discrete diffusion" OR "masked diffusion" OR "absorbing diffusion" OR "absorbing state" OR "discrete flow matching" OR "diffusion LLM" OR "diffusion large language model*" OR "text diffusion" OR ("non-autoregressive" AND diffusion)) AND (language OR text OR token* OR "language model*")}\par\smallskip
\noindent\textbf{S2}, 2026-10-15 10:55 UTC, 5 records.\\
{\ttfamily\footnotesize ("unified multimodal" OR "unified understanding and generation" OR "understanding and generation" OR "any-to-any") AND (diffusion OR "masked generative" OR "discrete flow" OR "flow matching") AND (multimodal OR "vision-language" OR "text-to-image")}\par\smallskip
\noindent\textbf{S1'}, 2026-10-15 12:55 UTC, 230 records.\\
{\ttfamily\footnotesize ("diffusion language model*" OR "language diffusion model*" OR "diffusion llm*" OR "diffusion lm*" OR dllm* OR "diffusion-based language model*" OR "diffusion-based llm*" OR "diffusion large language model*" OR "diffusion-based large language model*" OR "text diffusion" OR "block diffusion" OR "absorbing diffusion" OR "diffusion mllm*" OR "diffusion vlm*" OR "diffusion vision-language model*" OR dmllm* OR dvlm*) OR ((("discrete diffusion" OR "masked diffusion" OR "discrete flow*" OR "discrete state-space*") AND (language OR text OR token* OR "language model*" OR llm* OR vocabulary OR autoregressive OR reasoning)) OR (("discrete diffusion" OR "masked diffusion" OR "discrete flow*" OR "discrete state-space*") AND NOT (protein* OR molecul* OR antibod* OR peptide* OR rna OR dna OR drug* OR crystal* OR chemi* OR genom* OR biolog* OR "graph generation" OR "traveling salesman" OR recommend*) AND NOT (image* OR video* OR 3d OR audio OR visual* OR pixel* OR photo* OR segmentation))) OR (("absorbing state" AND (language OR text OR token* OR "language model*")) OR ("any-order" AND (language OR text OR token* OR "language model*")) OR "non-autoregressive language model*" OR ("non-autoregressive" AND diffusion AND NOT (protein* OR molecul* OR antibod* OR peptide* OR rna OR dna OR drug* OR crystal* OR chemi* OR genom* OR biolog* OR "graph generation" OR "traveling salesman" OR recommend*))) OR (diffusion AND ("text generation" OR "language generation" OR "natural language generation" OR "sequence-to-sequence" OR seq2seq OR "discrete sequence*" OR "speculative decoding" OR "categorical data" OR "discrete data" OR paraphras* OR "textual diffusion" OR "text modeling" OR "named entity" OR sentiment OR summarization OR dialog* OR dialogue OR "machine translation" OR "text style" OR commonsense OR "topic model*" OR "triple extraction" OR "information extraction") AND NOT (protein* OR molecul* OR antibod* OR peptide* OR rna OR dna OR drug* OR crystal* OR chemi* OR genom* OR biolog* OR "graph generation" OR "traveling salesman" OR recommend*)) OR (diffusion AND ("language model*" OR llm*) AND NOT (protein* OR molecul* OR antibod* OR peptide* OR rna OR dna OR drug* OR crystal* OR chemi* OR genom* OR biolog* OR "graph generation" OR "traveling salesman" OR recommend*) AND NOT (image* OR video* OR 3d OR audio OR visual* OR pixel* OR photo* OR segmentation))}\par\smallskip
\noindent\textbf{S2'}, 2026-10-15 12:55 UTC, 7 records.\\
{\ttfamily\footnotesize (("unified multimodal" OR "unified multi-modal" OR "unified understanding and generation" OR "understanding and generation" OR "any-to-any" OR "mixed-modal*" OR "mixed-modality") AND (diffusion OR "masked generative" OR "discrete flow*" OR "flow matching" OR "rectified flow" OR "discrete denoising") AND (multimodal OR "multi-modal" OR "vision-language" OR "text-to-image" OR "image generation" OR mllm* OR "visual understanding" OR "image understanding")) OR (("vision-language-action" OR vla*) AND ("discrete diffusion" OR "discrete denoising" OR "masked diffusion" OR "diffusion-based" OR "unified diffusion" OR "diffusion language model*" OR dllm*))}\par\smallskip
\noindent\textbf{S1'}, 2026-10-15 13:36 UTC, 230 records.\\
{\ttfamily\footnotesize ("diffusion language model*" OR "language diffusion model*" OR "diffusion llm*" OR "diffusion lm*" OR dllm OR dllms OR "diffusion-based language model*" OR "diffusion-based llm*" OR "diffusion large language model*" OR "diffusion-based large language model*" OR "text diffusion" OR "block diffusion" OR "absorbing diffusion" OR "diffusion mllm*" OR "diffusion vlm*" OR "diffusion vision-language model*" OR dmllm OR dmllms OR dvlm OR dvlms) OR ((("discrete diffusion" OR "masked diffusion" OR "discrete flow*" OR "discrete state-space*") AND (language OR text OR token* OR "language model*" OR llm* OR vocabulary OR autoregressive OR reasoning)) OR (("discrete diffusion" OR "masked diffusion" OR "discrete flow*" OR "discrete state-space*") AND NOT (protein* OR molecul* OR antibod* OR peptide* OR rna OR dna OR drug* OR crystal* OR chemi* OR genom* OR biolog* OR "graph generation" OR "traveling salesman" OR recommend*) AND NOT (image* OR video* OR 3d OR audio OR visual* OR pixel* OR photo* OR segmentation))) OR (("absorbing state" AND (language OR text OR token* OR "language model*")) OR ("any-order" AND (language OR text OR token* OR "language model*")) OR "non-autoregressive language model*" OR ("non-autoregressive" AND diffusion AND NOT (protein* OR molecul* OR antibod* OR peptide* OR rna OR dna OR drug* OR crystal* OR chemi* OR genom* OR biolog* OR "graph generation" OR "traveling salesman" OR recommend*))) OR (diffusion AND ("text generation" OR "language generation" OR "natural language generation" OR "sequence-to-sequence" OR seq2seq OR "discrete sequence*" OR "speculative decoding" OR "categorical data" OR "discrete data" OR paraphras* OR "textual diffusion" OR "text modeling" OR "named entity" OR sentiment OR summarization OR dialog* OR dialogue OR "machine translation" OR "text style" OR commonsense OR "topic model*" OR "triple extraction" OR "information extraction") AND NOT (protein* OR molecul* OR antibod* OR peptide* OR rna OR dna OR drug* OR crystal* OR chemi* OR genom* OR biolog* OR "graph generation" OR "traveling salesman" OR recommend*)) OR (diffusion AND ("language model*" OR llm*) AND NOT (protein* OR molecul* OR antibod* OR peptide* OR rna OR dna OR drug* OR crystal* OR chemi* OR genom* OR biolog* OR "graph generation" OR "traveling salesman" OR recommend*) AND NOT (image* OR video* OR 3d OR audio OR visual* OR pixel* OR photo* OR segmentation))}\par\smallskip
\noindent\textbf{S2'}, 2026-10-15 13:36 UTC, 7 records.\\
{\ttfamily\footnotesize (("unified multimodal" OR "unified multi-modal" OR "unified understanding and generation" OR "understanding and generation" OR "any-to-any" OR "mixed-modal*" OR "mixed-modality") AND (diffusion OR "masked generative" OR "discrete flow*" OR "flow matching" OR "rectified flow" OR "discrete denoising") AND (multimodal OR "multi-modal" OR "vision-language" OR "text-to-image" OR "image generation" OR mllm OR mllms OR "visual understanding" OR "image understanding")) OR (("vision-language-action" OR vla OR vlas) AND ("discrete diffusion" OR "discrete denoising" OR "masked diffusion" OR "diffusion-based" OR "unified diffusion" OR "diffusion language model*" OR dllm OR dllms))}\par\smallskip
\paragraph{Semantic Scholar Graph API (paper/batch references)}
\noindent\textbf{snowball-backward}, 2026-10-15 10:56 UTC, 25509 records.\\
{\ttfamily\footnotesize seed set = 539 arXiv IDs of Awesome-DLMs@194aa6bd; one level}\par\smallskip
\paragraph{Semantic Scholar Graph API (paper/\{id\}/citations)}
\noindent\textbf{snowball-forward}, 2026-10-15 11:08 UTC, 26003 records.\\
{\ttfamily\footnotesize seed set = 539 arXiv IDs of Awesome-DLMs@194aa6bd; one level}\par\smallskip
\paragraph{arXiv API (export.arxiv.org), revised string, part A1}
\noindent\textbf{S1':A1}, 2026-10-15 12:55 UTC, 992 records.\\
{\ttfamily\footnotesize (((ti:"diffusion language model" OR abs:"diffusion language model") OR (ti:"diffusion language models" OR abs:"diffusion language models")) OR ((ti:"language diffusion model" OR abs:"language diffusion model") OR (ti:"language diffusion models" OR abs:"language diffusion models")) OR ((ti:"diffusion llm" OR abs:"diffusion llm") OR (ti:"diffusion llms" OR abs:"diffusion llms")) OR ((ti:"diffusion lm" OR abs:"diffusion lm") OR (ti:"diffusion lms" OR abs:"diffusion lms")) OR (ti:dllm* OR abs:dllm*) OR ((ti:"diffusion-based language model" OR abs:"diffusion-based language model") OR (ti:"diffusion-based language models" OR abs:"diffusion-based language models")) OR ((ti:"diffusion-based llm" OR abs:"diffusion-based llm") OR (ti:"diffusion-based llms" OR abs:"diffusion-based llms")) OR ((ti:"diffusion large language model" OR abs:"diffusion large language model") OR (ti:"diffusion large language models" OR abs:"diffusion large language models")) OR ((ti:"diffusion-based large language model" OR abs:"diffusion-based large language model") OR (ti:"diffusion-based large language models" OR abs:"diffusion-based large language models")) OR (ti:"text diffusion" OR abs:"text diffusion") OR (ti:"block diffusion" OR abs:"block diffusion") OR (ti:"absorbing diffusion" OR abs:"absorbing diffusion") OR ((ti:"diffusion mllm" OR abs:"diffusion mllm") OR (ti:"diffusion mllms" OR abs:"diffusion mllms")) OR ((ti:"diffusion vlm" OR abs:"diffusion vlm") OR (ti:"diffusion vlms" OR abs:"diffusion vlms")) OR ((ti:"diffusion vision-language model" OR abs:"diffusion vision-language model") OR (ti:"diffusion vision-language models" OR abs:"diffusion vision-language models")) OR (ti:dmllm* OR abs:dmllm*) OR (ti:dvlm* OR abs:dvlm*)) AND (cat:cs.CL OR cat:cs.LG OR cat:cs.CV OR cat:cs.AI OR cat:stat.ML) AND submittedDate:[202101010000 TO 202610152359]}\par\smallskip
\noindent\textbf{S1':A1}, 2026-10-15 13:18 UTC, 991 records.\\
{\ttfamily\footnotesize (((ti:"diffusion language model" OR abs:"diffusion language model") OR (ti:"diffusion language models" OR abs:"diffusion language models")) OR ((ti:"language diffusion model" OR abs:"language diffusion model") OR (ti:"language diffusion models" OR abs:"language diffusion models")) OR ((ti:"diffusion llm" OR abs:"diffusion llm") OR (ti:"diffusion llms" OR abs:"diffusion llms")) OR ((ti:"diffusion lm" OR abs:"diffusion lm") OR (ti:"diffusion lms" OR abs:"diffusion lms")) OR (ti:dllm OR abs:dllm) OR ((ti:"diffusion-based language model" OR abs:"diffusion-based language model") OR (ti:"diffusion-based language models" OR abs:"diffusion-based language models")) OR ((ti:"diffusion-based llm" OR abs:"diffusion-based llm") OR (ti:"diffusion-based llms" OR abs:"diffusion-based llms")) OR ((ti:"diffusion large language model" OR abs:"diffusion large language model") OR (ti:"diffusion large language models" OR abs:"diffusion large language models")) OR ((ti:"diffusion-based large language model" OR abs:"diffusion-based large language model") OR (ti:"diffusion-based large language models" OR abs:"diffusion-based large language models")) OR (ti:"text diffusion" OR abs:"text diffusion") OR (ti:"block diffusion" OR abs:"block diffusion") OR (ti:"absorbing diffusion" OR abs:"absorbing diffusion") OR ((ti:"diffusion mllm" OR abs:"diffusion mllm") OR (ti:"diffusion mllms" OR abs:"diffusion mllms")) OR ((ti:"diffusion vlm" OR abs:"diffusion vlm") OR (ti:"diffusion vlms" OR abs:"diffusion vlms")) OR ((ti:"diffusion vision-language model" OR abs:"diffusion vision-language model") OR (ti:"diffusion vision-language models" OR abs:"diffusion vision-language models")) OR (ti:dmllm OR abs:dmllm) OR (ti:dvlm OR abs:dvlm)) AND (cat:cs.CL OR cat:cs.LG OR cat:cs.CV OR cat:cs.AI OR cat:stat.ML) AND submittedDate:[202101010000 TO 202610152359]}\par\smallskip
\noindent\textbf{S1':A1}, 2026-10-15 13:36 UTC, 991 records.\\
{\ttfamily\footnotesize (((ti:"diffusion language model" OR abs:"diffusion language model") OR (ti:"diffusion language models" OR abs:"diffusion language models")) OR ((ti:"language diffusion model" OR abs:"language diffusion model") OR (ti:"language diffusion models" OR abs:"language diffusion models")) OR ((ti:"diffusion llm" OR abs:"diffusion llm") OR (ti:"diffusion llms" OR abs:"diffusion llms")) OR ((ti:"diffusion lm" OR abs:"diffusion lm") OR (ti:"diffusion lms" OR abs:"diffusion lms")) OR (ti:dllm OR abs:dllm) OR (ti:dllms OR abs:dllms) OR ((ti:"diffusion-based language model" OR abs:"diffusion-based language model") OR (ti:"diffusion-based language models" OR abs:"diffusion-based language models")) OR ((ti:"diffusion-based llm" OR abs:"diffusion-based llm") OR (ti:"diffusion-based llms" OR abs:"diffusion-based llms")) OR ((ti:"diffusion large language model" OR abs:"diffusion large language model") OR (ti:"diffusion large language models" OR abs:"diffusion large language models")) OR ((ti:"diffusion-based large language model" OR abs:"diffusion-based large language model") OR (ti:"diffusion-based large language models" OR abs:"diffusion-based large language models")) OR (ti:"text diffusion" OR abs:"text diffusion") OR (ti:"block diffusion" OR abs:"block diffusion") OR (ti:"absorbing diffusion" OR abs:"absorbing diffusion") OR ((ti:"diffusion mllm" OR abs:"diffusion mllm") OR (ti:"diffusion mllms" OR abs:"diffusion mllms")) OR ((ti:"diffusion vlm" OR abs:"diffusion vlm") OR (ti:"diffusion vlms" OR abs:"diffusion vlms")) OR ((ti:"diffusion vision-language model" OR abs:"diffusion vision-language model") OR (ti:"diffusion vision-language models" OR abs:"diffusion vision-language models")) OR (ti:dmllm OR abs:dmllm) OR (ti:dmllms OR abs:dmllms) OR (ti:dvlm OR abs:dvlm) OR (ti:dvlms OR abs:dvlms)) AND (cat:cs.CL OR cat:cs.LG OR cat:cs.CV OR cat:cs.AI OR cat:stat.ML) AND submittedDate:[202101010000 TO 202610152359]}\par\smallskip
\paragraph{arXiv API (export.arxiv.org), revised string, part A2}
\noindent\textbf{S1':A2}, 2026-10-15 12:55 UTC, 860 records.\\
{\ttfamily\footnotesize ((((ti:"discrete diffusion" OR abs:"discrete diffusion") OR (ti:"masked diffusion" OR abs:"masked diffusion") OR ((ti:"discrete flow" OR abs:"discrete flow") OR (ti:"discrete flows" OR abs:"discrete flows")) OR ((ti:"discrete state-space" OR abs:"discrete state-space") OR (ti:"discrete state-spaces" OR abs:"discrete state-spaces"))) AND ((ti:language OR abs:language) OR (ti:text OR abs:text) OR (ti:token* OR abs:token*) OR ((ti:"language model" OR abs:"language model") OR (ti:"language models" OR abs:"language models")) OR (ti:llm* OR abs:llm*) OR (ti:vocabulary OR abs:vocabulary) OR (ti:autoregressive OR abs:autoregressive) OR (ti:reasoning OR abs:reasoning))) OR (((((ti:"discrete diffusion" OR abs:"discrete diffusion") OR (ti:"masked diffusion" OR abs:"masked diffusion") OR ((ti:"discrete flow" OR abs:"discrete flow") OR (ti:"discrete flows" OR abs:"discrete flows")) OR ((ti:"discrete state-space" OR abs:"discrete state-space") OR (ti:"discrete state-spaces" OR abs:"discrete state-spaces")))) ANDNOT ((ti:protein* OR abs:protein*) OR (ti:molecul* OR abs:molecul*) OR (ti:antibod* OR abs:antibod*) OR (ti:peptide* OR abs:peptide*) OR (ti:rna OR abs:rna) OR (ti:dna OR abs:dna) OR (ti:drug* OR abs:drug*) OR (ti:crystal* OR abs:crystal*) OR (ti:chemi* OR abs:chemi*) OR (ti:genom* OR abs:genom*) OR (ti:biolog* OR abs:biolog*) OR (ti:"graph generation" OR abs:"graph generation") OR (ti:"traveling salesman" OR abs:"traveling salesman") OR (ti:recommend* OR abs:recommend*))) ANDNOT ((ti:image* OR abs:image*) OR (ti:video* OR abs:video*) OR (ti:3d OR abs:3d) OR (ti:audio OR abs:audio) OR (ti:visual* OR abs:visual*) OR (ti:pixel* OR abs:pixel*) OR (ti:photo* OR abs:photo*) OR (ti:segmentation OR abs:segmentation)))) AND (cat:cs.CL OR cat:cs.LG OR cat:cs.CV OR cat:cs.AI OR cat:stat.ML) AND submittedDate:[202101010000 TO 202610152359]}\par\smallskip
\noindent\textbf{S1':A2}, 2026-10-15 13:18 UTC, 837 records.\\
{\ttfamily\footnotesize ((((ti:"discrete diffusion" OR abs:"discrete diffusion") OR (ti:"masked diffusion" OR abs:"masked diffusion") OR ((ti:"discrete flow" OR abs:"discrete flow") OR (ti:"discrete flows" OR abs:"discrete flows")) OR ((ti:"discrete state-space" OR abs:"discrete state-space") OR (ti:"discrete state-spaces" OR abs:"discrete state-spaces"))) AND ((ti:language OR abs:language) OR (ti:text OR abs:text) OR ((ti:token OR abs:token) OR (ti:tokens OR abs:tokens) OR (ti:tokenizer OR abs:tokenizer)) OR ((ti:"language model" OR abs:"language model") OR (ti:"language models" OR abs:"language models")) OR (ti:llm OR abs:llm) OR (ti:vocabulary OR abs:vocabulary) OR (ti:autoregressive OR abs:autoregressive) OR (ti:reasoning OR abs:reasoning))) OR (((((ti:"discrete diffusion" OR abs:"discrete diffusion") OR (ti:"masked diffusion" OR abs:"masked diffusion") OR ((ti:"discrete flow" OR abs:"discrete flow") OR (ti:"discrete flows" OR abs:"discrete flows")) OR ((ti:"discrete state-space" OR abs:"discrete state-space") OR (ti:"discrete state-spaces" OR abs:"discrete state-spaces")))) ANDNOT ((ti:protein OR abs:protein) OR (ti:molecule OR abs:molecule) OR (ti:molecular OR abs:molecular) OR (ti:antibody OR abs:antibody) OR (ti:peptide OR abs:peptide) OR (ti:rna OR abs:rna) OR (ti:dna OR abs:dna) OR (ti:drug OR abs:drug) OR (ti:crystal OR abs:crystal) OR (ti:chemical OR abs:chemical) OR (ti:chemistry OR abs:chemistry) OR (ti:genome OR abs:genome) OR (ti:genomic OR abs:genomic) OR (ti:biology OR abs:biology) OR (ti:biological OR abs:biological) OR (ti:"graph generation" OR abs:"graph generation") OR (ti:"traveling salesman" OR abs:"traveling salesman") OR (ti:recommendation OR abs:recommendation) OR (ti:recommender OR abs:recommender))) ANDNOT ((ti:image OR abs:image) OR (ti:video OR abs:video) OR (ti:3d OR abs:3d) OR (ti:audio OR abs:audio) OR (ti:visual OR abs:visual) OR (ti:visualization OR abs:visualization) OR (ti:pixel OR abs:pixel) OR (ti:photo OR abs:photo) OR (ti:photograph OR abs:photograph) OR (ti:segmentation OR abs:segmentation)))) AND (cat:cs.CL OR cat:cs.LG OR cat:cs.CV OR cat:cs.AI OR cat:stat.ML) AND submittedDate:[202101010000 TO 202610152359]}\par\smallskip
\paragraph{arXiv API (export.arxiv.org), revised string, part A3}
\noindent\textbf{S1':A3}, 2026-10-15 12:56 UTC, 225 records.\\
{\ttfamily\footnotesize (((ti:"absorbing state" OR abs:"absorbing state") AND ((ti:language OR abs:language) OR (ti:text OR abs:text) OR (ti:token* OR abs:token*) OR ((ti:"language model" OR abs:"language model") OR (ti:"language models" OR abs:"language models")))) OR ((ti:"any-order" OR abs:"any-order") AND ((ti:language OR abs:language) OR (ti:text OR abs:text) OR (ti:token* OR abs:token*) OR ((ti:"language model" OR abs:"language model") OR (ti:"language models" OR abs:"language models")))) OR ((ti:"non-autoregressive language model" OR abs:"non-autoregressive language model") OR (ti:"non-autoregressive language models" OR abs:"non-autoregressive language models")) OR (((ti:"non-autoregressive" OR abs:"non-autoregressive") AND (ti:diffusion OR abs:diffusion)) ANDNOT ((ti:protein* OR abs:protein*) OR (ti:molecul* OR abs:molecul*) OR (ti:antibod* OR abs:antibod*) OR (ti:peptide* OR abs:peptide*) OR (ti:rna OR abs:rna) OR (ti:dna OR abs:dna) OR (ti:drug* OR abs:drug*) OR (ti:crystal* OR abs:crystal*) OR (ti:chemi* OR abs:chemi*) OR (ti:genom* OR abs:genom*) OR (ti:biolog* OR abs:biolog*) OR (ti:"graph generation" OR abs:"graph generation") OR (ti:"traveling salesman" OR abs:"traveling salesman") OR (ti:recommend* OR abs:recommend*)))) AND (cat:cs.CL OR cat:cs.LG OR cat:cs.CV OR cat:cs.AI OR cat:stat.ML) AND submittedDate:[202101010000 TO 202610152359]}\par\smallskip
\noindent\textbf{S1':A3}, 2026-10-15 13:18 UTC, 225 records.\\
{\ttfamily\footnotesize (((ti:"absorbing state" OR abs:"absorbing state") AND ((ti:language OR abs:language) OR (ti:text OR abs:text) OR ((ti:token OR abs:token) OR (ti:tokens OR abs:tokens) OR (ti:tokenizer OR abs:tokenizer)) OR ((ti:"language model" OR abs:"language model") OR (ti:"language models" OR abs:"language models")))) OR ((ti:"any-order" OR abs:"any-order") AND ((ti:language OR abs:language) OR (ti:text OR abs:text) OR ((ti:token OR abs:token) OR (ti:tokens OR abs:tokens) OR (ti:tokenizer OR abs:tokenizer)) OR ((ti:"language model" OR abs:"language model") OR (ti:"language models" OR abs:"language models")))) OR ((ti:"non-autoregressive language model" OR abs:"non-autoregressive language model") OR (ti:"non-autoregressive language models" OR abs:"non-autoregressive language models")) OR (((ti:"non-autoregressive" OR abs:"non-autoregressive") AND (ti:diffusion OR abs:diffusion)) ANDNOT ((ti:protein OR abs:protein) OR (ti:molecule OR abs:molecule) OR (ti:molecular OR abs:molecular) OR (ti:antibody OR abs:antibody) OR (ti:peptide OR abs:peptide) OR (ti:rna OR abs:rna) OR (ti:dna OR abs:dna) OR (ti:drug OR abs:drug) OR (ti:crystal OR abs:crystal) OR (ti:chemical OR abs:chemical) OR (ti:chemistry OR abs:chemistry) OR (ti:genome OR abs:genome) OR (ti:genomic OR abs:genomic) OR (ti:biology OR abs:biology) OR (ti:biological OR abs:biological) OR (ti:"graph generation" OR abs:"graph generation") OR (ti:"traveling salesman" OR abs:"traveling salesman") OR (ti:recommendation OR abs:recommendation) OR (ti:recommender OR abs:recommender)))) AND (cat:cs.CL OR cat:cs.LG OR cat:cs.CV OR cat:cs.AI OR cat:stat.ML) AND submittedDate:[202101010000 TO 202610152359]}\par\smallskip
\paragraph{arXiv API (export.arxiv.org), revised string, part B}
\noindent\textbf{S1':B}, 2026-10-15 12:56 UTC, 906 records.\\
{\ttfamily\footnotesize (((ti:diffusion OR abs:diffusion) AND ((ti:"text generation" OR abs:"text generation") OR (ti:"language generation" OR abs:"language generation") OR (ti:"natural language generation" OR abs:"natural language generation") OR (ti:"sequence-to-sequence" OR abs:"sequence-to-sequence") OR (ti:seq2seq OR abs:seq2seq) OR ((ti:"discrete sequence" OR abs:"discrete sequence") OR (ti:"discrete sequences" OR abs:"discrete sequences")) OR (ti:"speculative decoding" OR abs:"speculative decoding") OR (ti:"categorical data" OR abs:"categorical data") OR (ti:"discrete data" OR abs:"discrete data") OR (ti:paraphras* OR abs:paraphras*) OR (ti:"textual diffusion" OR abs:"textual diffusion") OR (ti:"text modeling" OR abs:"text modeling") OR (ti:"named entity" OR abs:"named entity") OR (ti:sentiment OR abs:sentiment) OR (ti:summarization OR abs:summarization) OR (ti:dialog* OR abs:dialog*) OR (ti:dialogue OR abs:dialogue) OR (ti:"machine translation" OR abs:"machine translation") OR (ti:"text style" OR abs:"text style") OR (ti:commonsense OR abs:commonsense) OR ((ti:"topic model" OR abs:"topic model") OR (ti:"topic models" OR abs:"topic models")) OR (ti:"triple extraction" OR abs:"triple extraction") OR (ti:"information extraction" OR abs:"information extraction"))) ANDNOT ((ti:protein* OR abs:protein*) OR (ti:molecul* OR abs:molecul*) OR (ti:antibod* OR abs:antibod*) OR (ti:peptide* OR abs:peptide*) OR (ti:rna OR abs:rna) OR (ti:dna OR abs:dna) OR (ti:drug* OR abs:drug*) OR (ti:crystal* OR abs:crystal*) OR (ti:chemi* OR abs:chemi*) OR (ti:genom* OR abs:genom*) OR (ti:biolog* OR abs:biolog*) OR (ti:"graph generation" OR abs:"graph generation") OR (ti:"traveling salesman" OR abs:"traveling salesman") OR (ti:recommend* OR abs:recommend*))) AND (cat:cs.CL OR cat:cs.LG OR cat:cs.CV OR cat:cs.AI OR cat:stat.ML) AND submittedDate:[202101010000 TO 202610152359]}\par\smallskip
\noindent\textbf{S1':B}, 2026-10-15 13:18 UTC, 906 records.\\
{\ttfamily\footnotesize (((ti:diffusion OR abs:diffusion) AND ((ti:"text generation" OR abs:"text generation") OR (ti:"language generation" OR abs:"language generation") OR (ti:"natural language generation" OR abs:"natural language generation") OR (ti:"sequence-to-sequence" OR abs:"sequence-to-sequence") OR (ti:seq2seq OR abs:seq2seq) OR ((ti:"discrete sequence" OR abs:"discrete sequence") OR (ti:"discrete sequences" OR abs:"discrete sequences")) OR (ti:"speculative decoding" OR abs:"speculative decoding") OR (ti:"categorical data" OR abs:"categorical data") OR (ti:"discrete data" OR abs:"discrete data") OR ((ti:paraphrase OR abs:paraphrase) OR (ti:paraphrasing OR abs:paraphrasing)) OR (ti:"textual diffusion" OR abs:"textual diffusion") OR (ti:"text modeling" OR abs:"text modeling") OR (ti:"named entity" OR abs:"named entity") OR (ti:sentiment OR abs:sentiment) OR (ti:summarization OR abs:summarization) OR ((ti:dialog OR abs:dialog) OR (ti:dialogue OR abs:dialogue)) OR (ti:dialogue OR abs:dialogue) OR (ti:"machine translation" OR abs:"machine translation") OR (ti:"text style" OR abs:"text style") OR (ti:commonsense OR abs:commonsense) OR ((ti:"topic model" OR abs:"topic model") OR (ti:"topic models" OR abs:"topic models")) OR (ti:"triple extraction" OR abs:"triple extraction") OR (ti:"information extraction" OR abs:"information extraction"))) ANDNOT ((ti:protein OR abs:protein) OR (ti:molecule OR abs:molecule) OR (ti:molecular OR abs:molecular) OR (ti:antibody OR abs:antibody) OR (ti:peptide OR abs:peptide) OR (ti:rna OR abs:rna) OR (ti:dna OR abs:dna) OR (ti:drug OR abs:drug) OR (ti:crystal OR abs:crystal) OR (ti:chemical OR abs:chemical) OR (ti:chemistry OR abs:chemistry) OR (ti:genome OR abs:genome) OR (ti:genomic OR abs:genomic) OR (ti:biology OR abs:biology) OR (ti:biological OR abs:biological) OR (ti:"graph generation" OR abs:"graph generation") OR (ti:"traveling salesman" OR abs:"traveling salesman") OR (ti:recommendation OR abs:recommendation) OR (ti:recommender OR abs:recommender))) AND (cat:cs.CL OR cat:cs.LG OR cat:cs.CV OR cat:cs.AI OR cat:stat.ML) AND submittedDate:[202101010000 TO 202610152359]}\par\smallskip
\paragraph{arXiv API (export.arxiv.org), revised string, part C}
\noindent\textbf{S1':C}, 2026-10-15 12:56 UTC, 1786 records.\\
{\ttfamily\footnotesize ((((ti:diffusion OR abs:diffusion) AND (((ti:"language model" OR abs:"language model") OR (ti:"language models" OR abs:"language models")) OR (ti:llm* OR abs:llm*))) ANDNOT ((ti:protein* OR abs:protein*) OR (ti:molecul* OR abs:molecul*) OR (ti:antibod* OR abs:antibod*) OR (ti:peptide* OR abs:peptide*) OR (ti:rna OR abs:rna) OR (ti:dna OR abs:dna) OR (ti:drug* OR abs:drug*) OR (ti:crystal* OR abs:crystal*) OR (ti:chemi* OR abs:chemi*) OR (ti:genom* OR abs:genom*) OR (ti:biolog* OR abs:biolog*) OR (ti:"graph generation" OR abs:"graph generation") OR (ti:"traveling salesman" OR abs:"traveling salesman") OR (ti:recommend* OR abs:recommend*))) ANDNOT ((ti:image* OR abs:image*) OR (ti:video* OR abs:video*) OR (ti:3d OR abs:3d) OR (ti:audio OR abs:audio) OR (ti:visual* OR abs:visual*) OR (ti:pixel* OR abs:pixel*) OR (ti:photo* OR abs:photo*) OR (ti:segmentation OR abs:segmentation))) AND (cat:cs.CL OR cat:cs.LG OR cat:cs.CV OR cat:cs.AI OR cat:stat.ML) AND submittedDate:[202101010000 TO 202610152359]}\par\smallskip
\noindent\textbf{S1':C}, 2026-10-15 13:19 UTC, 1359 records.\\
{\ttfamily\footnotesize ((((ti:diffusion OR abs:diffusion) AND (((ti:"language model" OR abs:"language model") OR (ti:"language models" OR abs:"language models")) OR (ti:llm OR abs:llm))) ANDNOT ((ti:protein OR abs:protein) OR (ti:molecule OR abs:molecule) OR (ti:molecular OR abs:molecular) OR (ti:antibody OR abs:antibody) OR (ti:peptide OR abs:peptide) OR (ti:rna OR abs:rna) OR (ti:dna OR abs:dna) OR (ti:drug OR abs:drug) OR (ti:crystal OR abs:crystal) OR (ti:chemical OR abs:chemical) OR (ti:chemistry OR abs:chemistry) OR (ti:genome OR abs:genome) OR (ti:genomic OR abs:genomic) OR (ti:biology OR abs:biology) OR (ti:biological OR abs:biological) OR (ti:"graph generation" OR abs:"graph generation") OR (ti:"traveling salesman" OR abs:"traveling salesman") OR (ti:recommendation OR abs:recommendation) OR (ti:recommender OR abs:recommender))) ANDNOT ((ti:image OR abs:image) OR (ti:video OR abs:video) OR (ti:3d OR abs:3d) OR (ti:audio OR abs:audio) OR (ti:visual OR abs:visual) OR (ti:visualization OR abs:visualization) OR (ti:pixel OR abs:pixel) OR (ti:photo OR abs:photo) OR (ti:photograph OR abs:photograph) OR (ti:segmentation OR abs:segmentation))) AND (cat:cs.CL OR cat:cs.LG OR cat:cs.CV OR cat:cs.AI OR cat:stat.ML) AND submittedDate:[202101010000 TO 202610152359]}\par\smallskip
\paragraph{arXiv API (export.arxiv.org), revised string, part M}
\noindent\textbf{S2':M}, 2026-10-15 12:57 UTC, 240 records.\\
{\ttfamily\footnotesize (((ti:"unified multimodal" OR abs:"unified multimodal") OR (ti:"unified multi-modal" OR abs:"unified multi-modal") OR (ti:"unified understanding and generation" OR abs:"unified understanding and generation") OR (ti:"understanding and generation" OR abs:"understanding and generation") OR (ti:"any-to-any" OR abs:"any-to-any") OR (ti:mixed-modal* OR abs:mixed-modal*) OR (ti:"mixed-modality" OR abs:"mixed-modality")) AND ((ti:diffusion OR abs:diffusion) OR (ti:"masked generative" OR abs:"masked generative") OR ((ti:"discrete flow" OR abs:"discrete flow") OR (ti:"discrete flows" OR abs:"discrete flows")) OR (ti:"flow matching" OR abs:"flow matching") OR (ti:"rectified flow" OR abs:"rectified flow") OR (ti:"discrete denoising" OR abs:"discrete denoising")) AND ((ti:multimodal OR abs:multimodal) OR (ti:"multi-modal" OR abs:"multi-modal") OR (ti:"vision-language" OR abs:"vision-language") OR (ti:"text-to-image" OR abs:"text-to-image") OR (ti:"image generation" OR abs:"image generation") OR (ti:mllm* OR abs:mllm*) OR (ti:"visual understanding" OR abs:"visual understanding") OR (ti:"image understanding" OR abs:"image understanding"))) AND (cat:cs.CL OR cat:cs.LG OR cat:cs.CV OR cat:cs.AI OR cat:stat.ML) AND submittedDate:[202101010000 TO 202610152359]}\par\smallskip
\noindent\textbf{S2':M}, 2026-10-15 13:19 UTC, 240 records.\\
{\ttfamily\footnotesize (((ti:"unified multimodal" OR abs:"unified multimodal") OR (ti:"unified multi-modal" OR abs:"unified multi-modal") OR (ti:"unified understanding and generation" OR abs:"unified understanding and generation") OR (ti:"understanding and generation" OR abs:"understanding and generation") OR (ti:"any-to-any" OR abs:"any-to-any") OR ((ti:"mixed-modal" OR abs:"mixed-modal") OR (ti:"mixed-modals" OR abs:"mixed-modals")) OR (ti:"mixed-modality" OR abs:"mixed-modality")) AND ((ti:diffusion OR abs:diffusion) OR (ti:"masked generative" OR abs:"masked generative") OR ((ti:"discrete flow" OR abs:"discrete flow") OR (ti:"discrete flows" OR abs:"discrete flows")) OR (ti:"flow matching" OR abs:"flow matching") OR (ti:"rectified flow" OR abs:"rectified flow") OR (ti:"discrete denoising" OR abs:"discrete denoising")) AND ((ti:multimodal OR abs:multimodal) OR (ti:"multi-modal" OR abs:"multi-modal") OR (ti:"vision-language" OR abs:"vision-language") OR (ti:"text-to-image" OR abs:"text-to-image") OR (ti:"image generation" OR abs:"image generation") OR (ti:mllm OR abs:mllm) OR (ti:"visual understanding" OR abs:"visual understanding") OR (ti:"image understanding" OR abs:"image understanding"))) AND (cat:cs.CL OR cat:cs.LG OR cat:cs.CV OR cat:cs.AI OR cat:stat.ML) AND submittedDate:[202101010000 TO 202610152359]}\par\smallskip
\noindent\textbf{S2':M}, 2026-10-15 13:36 UTC, 240 records.\\
{\ttfamily\footnotesize (((ti:"unified multimodal" OR abs:"unified multimodal") OR (ti:"unified multi-modal" OR abs:"unified multi-modal") OR (ti:"unified understanding and generation" OR abs:"unified understanding and generation") OR (ti:"understanding and generation" OR abs:"understanding and generation") OR (ti:"any-to-any" OR abs:"any-to-any") OR ((ti:"mixed-modal" OR abs:"mixed-modal") OR (ti:"mixed-modals" OR abs:"mixed-modals")) OR (ti:"mixed-modality" OR abs:"mixed-modality")) AND ((ti:diffusion OR abs:diffusion) OR (ti:"masked generative" OR abs:"masked generative") OR ((ti:"discrete flow" OR abs:"discrete flow") OR (ti:"discrete flows" OR abs:"discrete flows")) OR (ti:"flow matching" OR abs:"flow matching") OR (ti:"rectified flow" OR abs:"rectified flow") OR (ti:"discrete denoising" OR abs:"discrete denoising")) AND ((ti:multimodal OR abs:multimodal) OR (ti:"multi-modal" OR abs:"multi-modal") OR (ti:"vision-language" OR abs:"vision-language") OR (ti:"text-to-image" OR abs:"text-to-image") OR (ti:"image generation" OR abs:"image generation") OR (ti:mllm OR abs:mllm) OR (ti:mllms OR abs:mllms) OR (ti:"visual understanding" OR abs:"visual understanding") OR (ti:"image understanding" OR abs:"image understanding"))) AND (cat:cs.CL OR cat:cs.LG OR cat:cs.CV OR cat:cs.AI OR cat:stat.ML) AND submittedDate:[202101010000 TO 202610152359]}\par\smallskip
\paragraph{arXiv API (export.arxiv.org), revised string, part VLA}
\noindent\textbf{S2':VLA}, 2026-10-15 12:57 UTC, 74 records.\\
{\ttfamily\footnotesize (((ti:"vision-language-action" OR abs:"vision-language-action") OR (ti:vla* OR abs:vla*)) AND ((ti:"discrete diffusion" OR abs:"discrete diffusion") OR (ti:"discrete denoising" OR abs:"discrete denoising") OR (ti:"masked diffusion" OR abs:"masked diffusion") OR (ti:"diffusion-based" OR abs:"diffusion-based") OR (ti:"unified diffusion" OR abs:"unified diffusion") OR ((ti:"diffusion language model" OR abs:"diffusion language model") OR (ti:"diffusion language models" OR abs:"diffusion language models")) OR (ti:dllm* OR abs:dllm*))) AND (cat:cs.CL OR cat:cs.LG OR cat:cs.CV OR cat:cs.AI OR cat:stat.ML) AND submittedDate:[202101010000 TO 202610152359]}\par\smallskip
\noindent\textbf{S2':VLA}, 2026-10-15 13:19 UTC, 71 records.\\
{\ttfamily\footnotesize (((ti:"vision-language-action" OR abs:"vision-language-action") OR (ti:vla OR abs:vla)) AND ((ti:"discrete diffusion" OR abs:"discrete diffusion") OR (ti:"discrete denoising" OR abs:"discrete denoising") OR (ti:"masked diffusion" OR abs:"masked diffusion") OR (ti:"diffusion-based" OR abs:"diffusion-based") OR (ti:"unified diffusion" OR abs:"unified diffusion") OR ((ti:"diffusion language model" OR abs:"diffusion language model") OR (ti:"diffusion language models" OR abs:"diffusion language models")) OR (ti:dllm OR abs:dllm))) AND (cat:cs.CL OR cat:cs.LG OR cat:cs.CV OR cat:cs.AI OR cat:stat.ML) AND submittedDate:[202101010000 TO 202610152359]}\par\smallskip
\noindent\textbf{S2':VLA}, 2026-10-15 13:36 UTC, 71 records.\\
{\ttfamily\footnotesize (((ti:"vision-language-action" OR abs:"vision-language-action") OR (ti:vla OR abs:vla) OR (ti:vlas OR abs:vlas)) AND ((ti:"discrete diffusion" OR abs:"discrete diffusion") OR (ti:"discrete denoising" OR abs:"discrete denoising") OR (ti:"masked diffusion" OR abs:"masked diffusion") OR (ti:"diffusion-based" OR abs:"diffusion-based") OR (ti:"unified diffusion" OR abs:"unified diffusion") OR ((ti:"diffusion language model" OR abs:"diffusion language model") OR (ti:"diffusion language models" OR abs:"diffusion language models")) OR (ti:dllm OR abs:dllm) OR (ti:dllms OR abs:dllms))) AND (cat:cs.CL OR cat:cs.LG OR cat:cs.CV OR cat:cs.AI OR cat:stat.ML) AND submittedDate:[202101010000 TO 202610152359]}\par\smallskip
\paragraph{Scopus (subscription export, ingest tools/ingest\_subscription.py)}
\noindent\textbf{Q1: union of S1, S2, S1' (A1, A2-r1, A3-r1..r3) and S2' (M, VLA), D0031 combined query}, 2026-10-15 16:50 UTC, 876 records.\\
{\ttfamily\footnotesize TITLE-ABS-KEY(("diffusion language model*" OR "language diffusion model*" OR "diffusion llm*" OR "diffusion lm*" OR dllm OR dllms OR "diffusion-based language model*" OR "diffusion-based llm*" OR "diffusion large language model*" OR "diffusion-based large language model*" OR "text diffusion" OR "block diffusion" OR "absorbing diffusion" OR "diffusion mllm*" OR "diffusion vlm*" OR "diffusion vision-language model*" OR dmllm OR dmllms OR dvlm OR dvlms OR "non-autoregressive language model*") OR (("discrete diffusion" OR "masked diffusion" OR "discrete flow*" OR "discrete state-space*" OR "absorbing state" OR "any-order" OR "discrete flow matching" OR "diffusion LLM" OR ("non-autoregressive" AND diffusion)) AND (language OR text OR token* OR "language model*" OR llm* OR vocabulary OR autoregressive OR reasoning)) OR (("unified multimodal" OR "unified multi-modal" OR "unified understanding and generation" OR "understanding and generation" OR "any-to-any" OR "mixed-modal*" OR "mixed-modality") AND (diffusion OR "masked generative" OR "discrete flow*" OR "flow matching" OR "rectified flow" OR "discrete denoising") AND (multimodal OR "multi-modal" OR "vision-language" OR "text-to-image" OR "image generation" OR mllm OR mllms OR "visual understanding" OR "image understanding")) OR (("vision-language-action" OR vla OR vlas) AND ("discrete diffusion" OR "discrete denoising" OR "masked diffusion" OR "diffusion-based" OR "unified diffusion" OR "diffusion language model*" OR dllm OR dllms))) AND PUBYEAR > 2020 AND PUBYEAR < 2027 AND LANGUAGE(english)}\par\smallskip
\noindent\textbf{S1':A2-r2}, 2026-10-15 16:51 UTC, 530 records.\\
{\ttfamily\footnotesize verbatim part S1p-A2-r2 of research/review/corpus/subscription-search-sheet.md (Scopus section); query text not repeated in research/sources/search-raw/scopus/scopus-queries-20261016.md}\par\smallskip
\noindent\textbf{S1':A3-r4}, 2026-10-15 16:52 UTC, 134 records.\\
{\ttfamily\footnotesize verbatim part S1p-A3-r4 of research/review/corpus/subscription-search-sheet.md (Scopus section); query text not repeated in research/sources/search-raw/scopus/scopus-queries-20261016.md}\par\smallskip
\noindent\textbf{S1':B}, 2026-10-15 16:54 UTC, 1670 records.\\
{\ttfamily\footnotesize verbatim part S1p-B of research/review/corpus/subscription-search-sheet.md (Scopus section); query text not repeated in research/sources/search-raw/scopus/scopus-queries-20261016.md}\par\smallskip
\noindent\textbf{S1':C}, 2026-10-15 16:48 UTC, 836 records.\\
{\ttfamily\footnotesize verbatim part S1p-C of research/review/corpus/subscription-search-sheet.md (Scopus section); query text not repeated in research/sources/search-raw/scopus/scopus-queries-20261016.md}\par\smallskip
\paragraph{Web of Science Core Collection (subscription export, ingest tools/ingest\_subscription.py)}
\noindent\textbf{union}, 2026-10-15 18:22 UTC, 1000 records (hits reported 2267).\\
{\ttfamily\footnotesize see research/sources/search-raw/web-of-science/wos-queries-20261015.md}\par\smallskip
\noindent\textbf{union}, 2026-10-15 18:23 UTC, 1000 records (hits reported 2267).\\
{\ttfamily\footnotesize see research/sources/search-raw/web-of-science/wos-queries-20261015.md}\par\smallskip
\noindent\textbf{union}, 2026-10-15 18:24 UTC, 267 records (hits reported 2267).\\
{\ttfamily\footnotesize see research/sources/search-raw/web-of-science/wos-queries-20261015.md}\par\smallskip
\paragraph{IEEE Xplore (subscription export, ingest tools/ingest\_subscription.py)}
\noindent\textbf{I1: S1':B + S1':C (NOT dropped): diffusion AND (C + B word lists)}, 2026-10-15 18:27 UTC, 230 records (hits reported 1203).\\
{\ttfamily\footnotesize ("Document Title":diffusion OR "Abstract":diffusion OR "Author Keywords":diffusion) AND (("Document Title":"language model" OR "Abstract":"language model" OR "Author Keywords":"language model") OR ("Document Title":"language models" OR "Abstract":"language models" OR "Author Keywords":"language models") OR ("Document Title":llm* OR "Abstract":llm* OR "Author Keywords":llm*) OR ("Document Title":"text generation" OR "Abstract":"text generation" OR "Author Keywords":"text generation") OR ("Document Title":"language generation" OR "Abstract":"language generation" OR "Author Keywords":"language generation") OR ("Document Title":"sequence-to-sequence" OR "Abstract":"sequence-to-sequence" OR "Author Keywords":"sequence-to-sequence") OR ("Document Title":seq2seq OR "Abstract":seq2seq OR "Author Keywords":seq2seq) OR ("Document Title":"discrete sequence" OR "Abstract":"discrete sequence" OR "Author Keywords":"discrete sequence") OR ("Document Title":"discrete sequences" OR "Abstract":"discrete sequences" OR "Author Keywords":"discrete sequences") OR ("Document Title":"speculative decoding" OR "Abstract":"speculative decoding" OR "Author Keywords":"speculative decoding") OR ("Document Title":"categorical data" OR "Abstract":"categorical data" OR "Author Keywords":"categorical data") OR ("Document Title":"discrete data" OR "Abstract":"discrete data" OR "Author Keywords":"discrete data") OR ("Document Title":paraphras* OR "Abstract":paraphras* OR "Author Keywords":paraphras*) OR ("Document Title":"textual diffusion" OR "Abstract":"textual diffusion" OR "Author Keywords":"textual diffusion") OR ("Document Title":"text modeling" OR "Abstract":"text modeling" OR "Author Keywords":"text modeling") OR ("Document Title":"named entity" OR "Abstract":"named entity" OR "Author Keywords":"named entity") OR ("Document Title":sentiment OR "Abstract":sentiment OR "Author Keywords":sentiment) OR ("Document Title":summarization OR "Abstract":summarization OR "Author Keywords":summarization) OR ("Document Title":dialog* OR "Abstract":dialog* OR "Author Keywords":dialog*) OR ("Document Title":"machine translation" OR "Abstract":"machine translation" OR "Author Keywords":"machine translation") OR ("Document Title":"text style" OR "Abstract":"text style" OR "Author Keywords":"text style") OR ("Document Title":commonsense OR "Abstract":commonsense OR "Author Keywords":commonsense) OR ("Document Title":"topic model" OR "Abstract":"topic model" OR "Author Keywords":"topic model") OR ("Document Title":"topic models" OR "Abstract":"topic models" OR "Author Keywords":"topic models") OR ("Document Title":"triple extraction" OR "Abstract":"triple extraction" OR "Author Keywords":"triple extraction") OR ("Document Title":"information extraction" OR "Abstract":"information extraction" OR "Author Keywords":"information extraction"))}\par\smallskip
\noindent\textbf{I1: S1':B + S1':C (NOT dropped): diffusion AND (C + B word lists)}, 2026-10-15 18:38 UTC, 506 records (hits reported 1203).\\
{\ttfamily\footnotesize ("Document Title":diffusion OR "Abstract":diffusion OR "Author Keywords":diffusion) AND (("Document Title":"language model" OR "Abstract":"language model" OR "Author Keywords":"language model") OR ("Document Title":"language models" OR "Abstract":"language models" OR "Author Keywords":"language models") OR ("Document Title":llm* OR "Abstract":llm* OR "Author Keywords":llm*) OR ("Document Title":"text generation" OR "Abstract":"text generation" OR "Author Keywords":"text generation") OR ("Document Title":"language generation" OR "Abstract":"language generation" OR "Author Keywords":"language generation") OR ("Document Title":"sequence-to-sequence" OR "Abstract":"sequence-to-sequence" OR "Author Keywords":"sequence-to-sequence") OR ("Document Title":seq2seq OR "Abstract":seq2seq OR "Author Keywords":seq2seq) OR ("Document Title":"discrete sequence" OR "Abstract":"discrete sequence" OR "Author Keywords":"discrete sequence") OR ("Document Title":"discrete sequences" OR "Abstract":"discrete sequences" OR "Author Keywords":"discrete sequences") OR ("Document Title":"speculative decoding" OR "Abstract":"speculative decoding" OR "Author Keywords":"speculative decoding") OR ("Document Title":"categorical data" OR "Abstract":"categorical data" OR "Author Keywords":"categorical data") OR ("Document Title":"discrete data" OR "Abstract":"discrete data" OR "Author Keywords":"discrete data") OR ("Document Title":paraphras* OR "Abstract":paraphras* OR "Author Keywords":paraphras*) OR ("Document Title":"textual diffusion" OR "Abstract":"textual diffusion" OR "Author Keywords":"textual diffusion") OR ("Document Title":"text modeling" OR "Abstract":"text modeling" OR "Author Keywords":"text modeling") OR ("Document Title":"named entity" OR "Abstract":"named entity" OR "Author Keywords":"named entity") OR ("Document Title":sentiment OR "Abstract":sentiment OR "Author Keywords":sentiment) OR ("Document Title":summarization OR "Abstract":summarization OR "Author Keywords":summarization) OR ("Document Title":dialog* OR "Abstract":dialog* OR "Author Keywords":dialog*) OR ("Document Title":"machine translation" OR "Abstract":"machine translation" OR "Author Keywords":"machine translation") OR ("Document Title":"text style" OR "Abstract":"text style" OR "Author Keywords":"text style") OR ("Document Title":commonsense OR "Abstract":commonsense OR "Author Keywords":commonsense) OR ("Document Title":"topic model" OR "Abstract":"topic model" OR "Author Keywords":"topic model") OR ("Document Title":"topic models" OR "Abstract":"topic models" OR "Author Keywords":"topic models") OR ("Document Title":"triple extraction" OR "Abstract":"triple extraction" OR "Author Keywords":"triple extraction") OR ("Document Title":"information extraction" OR "Abstract":"information extraction" OR "Author Keywords":"information extraction"))}\par\smallskip
\noindent\textbf{I1: S1':B + S1':C (NOT dropped): diffusion AND (C + B word lists)}, 2026-10-15 18:41 UTC, 463 records (hits reported 1203).\\
{\ttfamily\footnotesize ("Document Title":diffusion OR "Abstract":diffusion OR "Author Keywords":diffusion) AND (("Document Title":"language model" OR "Abstract":"language model" OR "Author Keywords":"language model") OR ("Document Title":"language models" OR "Abstract":"language models" OR "Author Keywords":"language models") OR ("Document Title":llm* OR "Abstract":llm* OR "Author Keywords":llm*) OR ("Document Title":"text generation" OR "Abstract":"text generation" OR "Author Keywords":"text generation") OR ("Document Title":"language generation" OR "Abstract":"language generation" OR "Author Keywords":"language generation") OR ("Document Title":"sequence-to-sequence" OR "Abstract":"sequence-to-sequence" OR "Author Keywords":"sequence-to-sequence") OR ("Document Title":seq2seq OR "Abstract":seq2seq OR "Author Keywords":seq2seq) OR ("Document Title":"discrete sequence" OR "Abstract":"discrete sequence" OR "Author Keywords":"discrete sequence") OR ("Document Title":"discrete sequences" OR "Abstract":"discrete sequences" OR "Author Keywords":"discrete sequences") OR ("Document Title":"speculative decoding" OR "Abstract":"speculative decoding" OR "Author Keywords":"speculative decoding") OR ("Document Title":"categorical data" OR "Abstract":"categorical data" OR "Author Keywords":"categorical data") OR ("Document Title":"discrete data" OR "Abstract":"discrete data" OR "Author Keywords":"discrete data") OR ("Document Title":paraphras* OR "Abstract":paraphras* OR "Author Keywords":paraphras*) OR ("Document Title":"textual diffusion" OR "Abstract":"textual diffusion" OR "Author Keywords":"textual diffusion") OR ("Document Title":"text modeling" OR "Abstract":"text modeling" OR "Author Keywords":"text modeling") OR ("Document Title":"named entity" OR "Abstract":"named entity" OR "Author Keywords":"named entity") OR ("Document Title":sentiment OR "Abstract":sentiment OR "Author Keywords":sentiment) OR ("Document Title":summarization OR "Abstract":summarization OR "Author Keywords":summarization) OR ("Document Title":dialog* OR "Abstract":dialog* OR "Author Keywords":dialog*) OR ("Document Title":"machine translation" OR "Abstract":"machine translation" OR "Author Keywords":"machine translation") OR ("Document Title":"text style" OR "Abstract":"text style" OR "Author Keywords":"text style") OR ("Document Title":commonsense OR "Abstract":commonsense OR "Author Keywords":commonsense) OR ("Document Title":"topic model" OR "Abstract":"topic model" OR "Author Keywords":"topic model") OR ("Document Title":"topic models" OR "Abstract":"topic models" OR "Author Keywords":"topic models") OR ("Document Title":"triple extraction" OR "Abstract":"triple extraction" OR "Author Keywords":"triple extraction") OR ("Document Title":"information extraction" OR "Abstract":"information extraction" OR "Author Keywords":"information extraction"))}\par\smallskip
\noindent\textbf{I2: S1':A1 (all names)}, 2026-10-15 18:41 UTC, 101 records.\\
{\ttfamily\footnotesize (("Document Title":"diffusion language model" OR "Abstract":"diffusion language model" OR "Author Keywords":"diffusion language model") OR ("Document Title":"diffusion language models" OR "Abstract":"diffusion language models" OR "Author Keywords":"diffusion language models") OR ("Document Title":"language diffusion model" OR "Abstract":"language diffusion model" OR "Author Keywords":"language diffusion model") OR ("Document Title":"language diffusion models" OR "Abstract":"language diffusion models" OR "Author Keywords":"language diffusion models") OR ("Document Title":"diffusion llm" OR "Abstract":"diffusion llm" OR "Author Keywords":"diffusion llm") OR ("Document Title":"diffusion llms" OR "Abstract":"diffusion llms" OR "Author Keywords":"diffusion llms") OR ("Document Title":"diffusion lm" OR "Abstract":"diffusion lm" OR "Author Keywords":"diffusion lm") OR ("Document Title":"diffusion lms" OR "Abstract":"diffusion lms" OR "Author Keywords":"diffusion lms") OR ("Document Title":dllm OR "Abstract":dllm OR "Author Keywords":dllm) OR ("Document Title":dllms OR "Abstract":dllms OR "Author Keywords":dllms) OR ("Document Title":"diffusion-based language model" OR "Abstract":"diffusion-based language model" OR "Author Keywords":"diffusion-based language model") OR ("Document Title":"diffusion-based language models" OR "Abstract":"diffusion-based language models" OR "Author Keywords":"diffusion-based language models") OR ("Document Title":"diffusion-based llm" OR "Abstract":"diffusion-based llm" OR "Author Keywords":"diffusion-based llm") OR ("Document Title":"diffusion-based llms" OR "Abstract":"diffusion-based llms" OR "Author Keywords":"diffusion-based llms") OR ("Document Title":"diffusion large language model" OR "Abstract":"diffusion large language model" OR "Author Keywords":"diffusion large language model") OR ("Document Title":"diffusion large language models" OR "Abstract":"diffusion large language models" OR "Author Keywords":"diffusion large language models") OR ("Document Title":"diffusion-based large language model" OR "Abstract":"diffusion-based large language model" OR "Author Keywords":"diffusion-based large language model") OR ("Document Title":"diffusion-based large language models" OR "Abstract":"diffusion-based large language models" OR "Author Keywords":"diffusion-based large language models") OR ("Document Title":"text diffusion" OR "Abstract":"text diffusion" OR "Author Keywords":"text diffusion") OR ("Document Title":"block diffusion" OR "Abstract":"block diffusion" OR "Author Keywords":"block diffusion") OR ("Document Title":"absorbing diffusion" OR "Abstract":"absorbing diffusion" OR "Author Keywords":"absorbing diffusion") OR ("Document Title":"diffusion mllm" OR "Abstract":"diffusion mllm" OR "Author Keywords":"diffusion mllm") OR ("Document Title":"diffusion mllms" OR "Abstract":"diffusion mllms" OR "Author Keywords":"diffusion mllms") OR ("Document Title":"diffusion vlm" OR "Abstract":"diffusion vlm" OR "Author Keywords":"diffusion vlm") OR ("Document Title":"diffusion vlms" OR "Abstract":"diffusion vlms" OR "Author Keywords":"diffusion vlms") OR ("Document Title":"diffusion vision-language model" OR "Abstract":"diffusion vision-language model" OR "Author Keywords":"diffusion vision-language model") OR ("Document Title":"diffusion vision-language models" OR "Abstract":"diffusion vision-language models" OR "Author Keywords":"diffusion vision-language models") OR ("Document Title":dmllm OR "Abstract":dmllm OR "Author Keywords":dmllm) OR ("Document Title":dmllms OR "Abstract":dmllms OR "Author Keywords":dmllms) OR ("Document Title":dvlm OR "Abstract":dvlm OR "Author Keywords":dvlm) OR ("Document Title":dvlms OR "Abstract":dvlms OR "Author Keywords":dvlms))}\par\smallskip
\paragraph{ACM Digital Library (subscription export, ingest tools/ingest\_subscription.py)}
\noindent\textbf{unspecified}, 2026-10-15 19:01 UTC, 329 records.\\
{\ttfamily\footnotesize see research/sources/search-raw/acm-dl/acm-queries-20261015.md}\par\smallskip
\paragraph{SpringerLink (subscription export, ingest tools/ingest\_subscription.py)}
\noindent\textbf{union}, 2026-10-16 04:07 UTC, 106 records.\\
{\ttfamily\footnotesize see research/sources/search-raw/springerlink/springer-queries-20261016.md}\par\smallskip
\noindent\textbf{S1}, 2026-10-16 04:44 UTC, 757 records.\\
{\ttfamily\footnotesize see research/sources/search-raw/springerlink/springer-queries-20261016.md}\par\smallskip
\noindent\textbf{S2:top1000}, 2026-10-16 04:45 UTC, 1000 records.\\
{\ttfamily\footnotesize see research/sources/search-raw/springerlink/springer-queries-20261016.md}\par\smallskip
\noindent\textbf{unspecified}, 2026-10-16 04:44 UTC, 659 records.\\
{\ttfamily\footnotesize see research/sources/search-raw/springerlink/springer-queries-20261016.md}\par\smallskip
\noindent\textbf{unspecified}, 2026-10-16 04:44 UTC, 92 records.\\
{\ttfamily\footnotesize see research/sources/search-raw/springerlink/springer-queries-20261016.md}\par\smallskip
\paragraph{ScienceDirect (subscription export, ingest tools/ingest\_subscription.py)}
\noindent\textbf{unspecified}, 2026-10-16 08:47 UTC, 10 records.\\
{\ttfamily\footnotesize see research/sources/search-raw/sciencedirect/sd-queries-20261016.md}\par\smallskip
\noindent\textbf{unspecified}, 2026-10-16 08:47 UTC, 2 records.\\
{\ttfamily\footnotesize see research/sources/search-raw/sciencedirect/sd-queries-20261016.md}\par\smallskip
\noindent\textbf{unspecified}, 2026-10-16 08:53 UTC, 100 records.\\
{\ttfamily\footnotesize see research/sources/search-raw/sciencedirect/sd-queries-20261016.md}\par\smallskip
\noindent\textbf{unspecified}, 2026-10-16 08:54 UTC, 37 records.\\
{\ttfamily\footnotesize see research/sources/search-raw/sciencedirect/sd-queries-20261016.md}\par\smallskip
\noindent\textbf{unspecified}, 2026-10-16 08:54 UTC, 100 records.\\
{\ttfamily\footnotesize see research/sources/search-raw/sciencedirect/sd-queries-20261016.md}\par\smallskip
\noindent\textbf{unspecified}, 2026-10-16 08:55 UTC, 100 records.\\
{\ttfamily\footnotesize see research/sources/search-raw/sciencedirect/sd-queries-20261016.md}\par\smallskip
\noindent\textbf{unspecified}, 2026-10-16 08:57 UTC, 52 records.\\
{\ttfamily\footnotesize see research/sources/search-raw/sciencedirect/sd-queries-20261016.md}\par\smallskip
\noindent\textbf{unspecified}, 2026-10-16 08:58 UTC, 100 records.\\
{\ttfamily\footnotesize see research/sources/search-raw/sciencedirect/sd-queries-20261016.md}\par\smallskip
\noindent\textbf{unspecified}, 2026-10-16 08:59 UTC, 100 records.\\
{\ttfamily\footnotesize see research/sources/search-raw/sciencedirect/sd-queries-20261016.md}\par\smallskip
\noindent\textbf{unspecified}, 2026-10-16 09:00 UTC, 9 records.\\
{\ttfamily\footnotesize see research/sources/search-raw/sciencedirect/sd-queries-20261016.md}\par\smallskip
\noindent\textbf{unspecified}, 2026-10-16 09:01 UTC, 1 records.\\
{\ttfamily\footnotesize see research/sources/search-raw/sciencedirect/sd-queries-20261016.md}\par\smallskip
\noindent\textbf{unspecified}, 2026-10-16 09:02 UTC, 2 records.\\
{\ttfamily\footnotesize see research/sources/search-raw/sciencedirect/sd-queries-20261016.md}\par\smallskip
\paragraph{OpenReview API v2 (api2.openreview.net/notes/search)}
group restricts to the venue family; date (cdate>=2021-01-01), the second clause of the string and the phrase test are applied locally by tools/merge\_new\_records.py; result window capped at 10000 by the API. Query form: \texttt{term=<phrase>\&source=forum\&group=<venue group>}. 204 runs, 5182 records before deduplication:
\begin{itemize}\setlength\itemsep{0pt}
\item \texttt{S1:"diffusion language model"}: ICLR.cc 144; NeurIPS.cc 31; ICML.cc 35; TMLR 3
\item \texttt{S1:"discrete diffusion"}: ICLR.cc 411; NeurIPS.cc 95; ICML.cc 142; TMLR 16
\item \texttt{S1:"masked diffusion"}: ICLR.cc 286; NeurIPS.cc 50; ICML.cc 84; TMLR 6
\item \texttt{S1:"absorbing diffusion"}: ICLR.cc 6; NeurIPS.cc 4; ICML.cc 1; TMLR 0
\item \texttt{S1:"absorbing state"}: ICLR.cc 9; NeurIPS.cc 5; ICML.cc 6; TMLR 1
\item \texttt{S1:"discrete flow matching"}: ICLR.cc 70; NeurIPS.cc 15; ICML.cc 22; TMLR 2
\item \texttt{S1:"diffusion LLM"}: ICLR.cc 30; NeurIPS.cc 5; ICML.cc 10; TMLR 0
\item \texttt{S1:"diffusion large language model"}: ICLR.cc 23; NeurIPS.cc 0; ICML.cc 6; TMLR 0
\item \texttt{S1:"text diffusion"}: ICLR.cc 17; NeurIPS.cc 8; ICML.cc 9; TMLR 0
\item \texttt{S1:"non-autoregressive"}: ICLR.cc 138; NeurIPS.cc 29; ICML.cc 38; TMLR 11
\item \texttt{S2:"unified multimodal"}: ICLR.cc 236; NeurIPS.cc 27; ICML.cc 33; TMLR 8
\item \texttt{S2:"unified understanding and generation"}: ICLR.cc 17; NeurIPS.cc 1; ICML.cc 0; TMLR 1
\item \texttt{S2:"understanding and generation"}: ICLR.cc 257; NeurIPS.cc 54; ICML.cc 33; TMLR 12
\item \texttt{S2:"any-to-any"}: ICLR.cc 20; NeurIPS.cc 6; ICML.cc 8; TMLR 0
\item \texttt{S1':"diffusion language model"}: ICLR.cc 144; NeurIPS.cc 31; ICML.cc 35; TMLR 3
\item \texttt{S1':"language diffusion model"}: ICLR.cc 3; NeurIPS.cc 3; ICML.cc 0; TMLR 0
\item \texttt{S1':"diffusion llm"}: ICLR.cc 30; NeurIPS.cc 5; ICML.cc 10; TMLR 0
\item \texttt{S1':"diffusion lm"}: ICLR.cc 6; NeurIPS.cc 2; ICML.cc 2; TMLR 1
\item \texttt{S1':"dllm"}: ICLR.cc 118; NeurIPS.cc 13; ICML.cc 29; TMLR 2
\item \texttt{S1':"diffusion-based language model"}: ICLR.cc 0; NeurIPS.cc 0; ICML.cc 0; TMLR 0
\item \texttt{S1':"diffusion large language model"}: ICLR.cc 23; NeurIPS.cc 0; ICML.cc 6; TMLR 0
\item \texttt{S1':"text diffusion"}: ICLR.cc 17; NeurIPS.cc 8; ICML.cc 9; TMLR 0
\item \texttt{S1':"block diffusion"}: ICLR.cc 65; NeurIPS.cc 5; ICML.cc 13; TMLR 2
\item \texttt{S1':"absorbing diffusion"}: ICLR.cc 6; NeurIPS.cc 4; ICML.cc 1; TMLR 0
\item \texttt{S1':"diffusion mllm"}: ICLR.cc 1; NeurIPS.cc 1; ICML.cc 1; TMLR 0
\item \texttt{S1':"diffusion vision-language model"}: ICLR.cc 5; NeurIPS.cc 0; ICML.cc 0; TMLR 0
\item \texttt{S1':"dmllm"}: ICLR.cc 4; NeurIPS.cc 0; ICML.cc 0; TMLR 1
\item \texttt{S1':"discrete diffusion"}: ICLR.cc 411; NeurIPS.cc 95; ICML.cc 142; TMLR 16
\item \texttt{S1':"masked diffusion"}: ICLR.cc 286; NeurIPS.cc 50; ICML.cc 84; TMLR 6
\item \texttt{S1':"discrete flow"}: ICLR.cc 66; NeurIPS.cc 14; ICML.cc 23; TMLR 2
\item \texttt{S1':"discrete state-space"}: ICLR.cc 14; NeurIPS.cc 8; ICML.cc 5; TMLR 2
\item \texttt{S1':"any-order"}: ICLR.cc 53; NeurIPS.cc 16; ICML.cc 21; TMLR 3
\item \texttt{S1':"non-autoregressive language model"}: ICLR.cc 0; NeurIPS.cc 0; ICML.cc 0; TMLR 0
\item \texttt{S1':"diffusion text generation"}: ICLR.cc 1; NeurIPS.cc 0; ICML.cc 0; TMLR 0
\item \texttt{S1':"diffusion speculative decoding"}: ICLR.cc 9; NeurIPS.cc 0; ICML.cc 1; TMLR 1
\item \texttt{S1':"diffusion machine translation"}: ICLR.cc 0; NeurIPS.cc 0; ICML.cc 0; TMLR 0
\item \texttt{S1':"diffusion summarization"}: ICLR.cc 0; NeurIPS.cc 0; ICML.cc 0; TMLR 0
\item \texttt{S1':"diffusion dialogue"}: ICLR.cc 0; NeurIPS.cc 0; ICML.cc 0; TMLR 0
\item \texttt{S1':"diffusion paraphrase"}: ICLR.cc 0; NeurIPS.cc 0; ICML.cc 0; TMLR 0
\item \texttt{S1':"diffusion language generation"}: ICLR.cc 2; NeurIPS.cc 0; ICML.cc 0; TMLR 0
\item \texttt{S1':"diffusion large language model reasoning"}: ICLR.cc 4; NeurIPS.cc 0; ICML.cc 0; TMLR 0
\item \texttt{S2':"unified multimodal"}: ICLR.cc 236; NeurIPS.cc 27; ICML.cc 33; TMLR 8
\item \texttt{S2':"unified multi-modal"}: ICLR.cc 21; NeurIPS.cc 4; ICML.cc 4; TMLR 0
\item \texttt{S2':"unified understanding and generation"}: ICLR.cc 17; NeurIPS.cc 1; ICML.cc 0; TMLR 1
\item \texttt{S2':"understanding and generation"}: ICLR.cc 257; NeurIPS.cc 54; ICML.cc 33; TMLR 12
\item \texttt{S2':"any-to-any"}: ICLR.cc 20; NeurIPS.cc 6; ICML.cc 8; TMLR 0
\item \texttt{S2':"mixed-modal"}: ICLR.cc 7; NeurIPS.cc 1; ICML.cc 1; TMLR 0
\item \texttt{S2':"vision-language-action diffusion"}: ICLR.cc 0; NeurIPS.cc 0; ICML.cc 0; TMLR 0
\item \texttt{S2':"multimodal discrete diffusion"}: ICLR.cc 2; NeurIPS.cc 0; ICML.cc 1; TMLR 0
\item \texttt{S2':"multimodal masked diffusion"}: ICLR.cc 0; NeurIPS.cc 0; ICML.cc 0; TMLR 1
\item \texttt{S2':"multimodal flow matching"}: ICLR.cc 1; NeurIPS.cc 0; ICML.cc 0; TMLR 1
\end{itemize}
\end{small}


\setcounter{section}{1}
\section{Screening: confusion table}
\label{sec:s-screen}

The per-record screening table is still to be assembled (see the note above). The table below is moved from Section~\Mref{sec:method-screening}, where its reading is given.

\begin{table}[t]
\caption{Title-and-abstract screening, 1453 records. Rows: screener A; columns: screener B. Three-way agreement 0.865 ($\kappa=0.673$); pass versus exclude agreement 0.917 ($\kappa=0.728$). Agreement is between two language-model screeners; it is not agreement with a human.}
\label{tab:screening}
\small
\begin{tabular}{@{}lrrrr@{}}
\toprule
 & B include & B unclear & B exclude & Total \\
\midrule
A include & 1013 & 13 & 45 & 1071 \\ % src: tools/screening_agreement.py confusion table
A unclear & 63 & 34 & 51 & 148 \\
A exclude & 6 & 18 & 210 & 234 \\
\midrule
Total & 1082 & 65 & 306 & 1453 \\ % src: screening_agreement.py "A: {...} B: {...}" lines
\bottomrule
\end{tabular}
\end{table}

\setcounter{section}{2}
\section{Census: tables moved from the main text}
\label{sec:s-census}

\subsection{Census method in full (moved from Section~\Mref{sec:census-method})}
\label{sec:s-census-method}

The census predates the search update and samples the seed corpus only; its eligibility judgements used the protocol's eligibility rules as amended on 2026-10-15 (decision D0020), not criteria v1.6. % src: census-result-v2.md ("Sample"); decision D0020
\paragraph{Pool, order and stopping rule.}
The census pool is every record that neither screener excluded and that has an arXiv identifier, 1106 records. The pool was put in a random order with a fixed seed (20261015) before any coding. Coders worked down the order, checked full-text eligibility, and stopped when the primary coder's eligible count reached 150; it reached 151 at rank 176, so the sample is ranks 1 to 176. That count treated the 15 papers Codex could not read as ineligible; under the final eligibility the 150th eligible paper falls at rank 157, so 19 more ranks were coded than intended. The sample is still the first 176 ranks of a random order, so the extra ranks do not bias the proportions. % src: census-result-v2.md ("Sample", commit 8dd2af3); audits/2026-10-15-gate-census-v2.md P2; src: tools/census_sample.py docstring and code (SEED = 20261015); coverage/census-result-v2.md ("Sample")
Full texts are the arXiv HTML rendering of the latest version on the coding date, converted to plain text; four of the 176 ranks (23, 33, 90, 144) have no full text and are left out of every count. % src: tools/census_sample.py; census-result-v2.md ("Ranks 23, 33, 90 and 144 have no full text")
Code repositories were not opened; the census records what the paper says, not what the released code does. % src: coverage/census-codebook-v1.md "Rules"

\paragraph{What is coded.}
Per paper the coder first judges eligibility under the protocol. For an eligible paper the coder identifies the headline result, defined as the first quantitative result for the paper's own main generative method in the abstract (skipping speed-only and likelihood-only numbers), or else the first row of the first main results table for the final method (field F1). % src: coverage/census-codebook-v1.md F1; census-codebook-v1.1.md A-2
If the paper reports no generated-sample metric, F1 is not applicable and the reporting fields are not coded; this applies to 14 eligible papers. One further paper shows samples only (F1 QL); by codebook rule B-28 it enters the value-rule count but not F3 to F6, so the value rule is counted on 153 papers and F3 to F6 on 152. % src: census-result-v2.md ("Result", first paragraph); analysis/census/summary-random.log (denominator lines); codebook B-10 and B-28
The value rule used for the headline (F2) is coded with letters: A argmax or greedy commit; T temperature other than one; P top-$p$ or top-$k$; S sample-then-filter (values are sampled and positions chosen by a score of the sampled value); G guidance; M sequence-level selection such as best-of-$N$; C confidence-ranked or threshold unmasking with the token rule unstated (added in v1.2); U plain sampling at temperature one with no truncation or guidance, coded only when the paper says so; and NS not stated. A suffix \texttt{+NS} marks a stated letter with the base token rule unstated. % src: coverage/census-codebook-v1.md F2; census-codebook-v1.2.md B-15 (C), census-result-v2.md ("15 carry +NS")
``Non-plain or unstated'' means any code other than U. The other fields are F3 (a diversity or entropy measure paired with the headline), F4 (the headline reported at two or more step settings, counting denoising steps, function evaluations or tokens per step; exactly two settings is coded ``partial''), F5 (the value rule stated as held fixed across the sweep) and F6 (guidance held fixed). % src: census-codebook-v1.md F3 to F6; v1.1 A-1, A-8, A-9
Conservative codes (NS, not stated) are used whenever a paper is ambiguous. % src: census-codebook-v1.md "Rules"

\paragraph{Codebook versions.}
Codebook v1 was written on 2026-10-14 and used for a first, purposive census of 47 models that is kept as a pilot and not used as evidence here. Amendments v1.1 (A-1 to A-10, 2026-10-14) governed the coding of the random sample. Codebook v1.2 (B-1 to B-28, 2026-10-15) was written by a separate adjudicator after both codings of the random sample and after agreement was computed, from the coders' own ambiguity reports, without opening the row files or the papers; it governs the adjudication, and agreement is reported under v1.1. % src: coverage/census-codebook-v1.md, -v1.1.md, -v1.2.md ("Disclosure"); coverage/census-result-v1.md; decision D0020
Rules in v1.2 that settle a gap and move a count are named where they matter (B-10, B-15, B-19, B-24). % src: decision D0020 item 3

\paragraph{Coders.}
Coder~A (Claude Sonnet) coded all 176 ranks, blind to coder~B. Coder~B was Codex (\texttt{gpt-6-luna}) for ranks 1 to 91, which gave 86 rows (ranks 19 and 79 failed), and then Claude Opus for ranks 92 to 176 and for ranks 19 and 79, because Codex reached its usage limit. This makes coder~A and part of coder~B models of the same family, so their agreement may overstate the agreement of independent human coders. % src: coverage/census-result-v2.md ("Process" step 1 and "Limits"); decision D0015; census/coder-B-run-log.jsonl (model gpt-6-luna in every row)
Codex could not read the paper for 15 ranks among 1 to 88; its codes there are tool-failure placeholders, and those rows were coded from the paper at adjudication. % src: census/adjudication-notes-r001-r088.md ("Coder B (Codex) could not read the file for 15 ranks"); census-result-v2.md
\todo{exact model identifiers and versions for coder~A and for coder~B (Claude Opus); the run logs record Codex only}

\paragraph{Agreement before reconciliation.}
Agreement between the two coders is computed on the 172 ranks with text before any reconciliation, under v1.1, and again without the 15 ranks Codex could not read; the 158 field-level disagreements are kept. % src: analysis/census/agreement-random.log and agreement-random-excl-readfail.log; census/disagreements-random.csv (158 rows)

\paragraph{Adjudication.}
Two adjudicators (Claude Opus; ranks 1 to 88 and 89 to 176) applied codebook v1.2 and checked the paper for each disputed cell, and one harmoniser then fixed six readings that the halves had settled differently. Adjudicators read both coders' rows and notes and were not blind. Of the 158 field-level disagreements, 73 were resolved to coder~A's value, 66 to coder~B's and 19 to neither. In addition, 45 cells on which both coders had agreed were reopened under a v1.2 rule, mostly because the coders' own notes showed its trigger. Three of them (ranks 114, 120, 141) were kept under the harmoniser's standard that an agreed value is changed when it cannot express a code that only v1.2 has (decision D0021). Of the 45, 42 affect a count, and the \emph{agreed cells restored} sensitivity of Section~\Mref{sec:census} undoes those 42. % src: census-result-v2.md (sensitivity paragraph, commit 8dd2af3); audits/2026-10-15-gate-census-v2.md C3, C8; decision D0021; adjudication counts: census/adjudication-notes-r001-r088.md (99 rows: A 59, B 31, neither 9); adjudication-notes-r089-r176.md (48 field rows: A 12, B 27, neither 9; 11 derived nonplain rows: A 2, B 8, neither 1); sum A 73, B 66, neither 19 = 158; census/harmonisation-notes.md (Date, harmoniser, six readings); census-result-v2.md ("45 cells")
The final rows, with both coders' values and the basis for every changed or disputed cell, are in \texttt{research/review/census/census-random-final.jsonl} and in the supplement (S3). % src: census-result-v2.md ("Process" step 4)

\paragraph{A blindness incident.}
Coder~B's Claude batch for ranks 92 to 128 left scratch copies of its rows and an append script in a shared scratch folder. The first batch of coder~A (ranks 1 to 40) ran that script once, so it may have seen coder~B's rows for ranks 92 to 128, which coder~A did not code. The copies were deleted and later batches ran as fresh agents in private scratch folders. We report it as a possible leak of one coder's labels into the other's workspace; we know of no effect on coder~A's codes. % src: decision D0019 item 2; census-result-v2.md ("Limits")



\subsection{Gap-filling rulings}

The codebooks and the 176-row final table are still to be assembled (see the note above). The table below is moved from Section~\Mref{sec:census-sensitivity}, where its reading is given.

\begin{table}[t]
\caption{Papers decided by each gap-filling ruling of codebook v1.2, among the value-rule denominator (153 papers). Columns give the value-rule class the ruling produced: non-plain, not stated (NS), plain stated (U). B-10 decided 15 ranks, of which 12 are in the value-rule count.}
\label{tab:census-gap}
\small
\begin{tabular}{@{}llrrrr@{}}
\toprule
Ruling & Subject & Papers & Non-plain & NS & U \\
\midrule
B-15 & code C, confidence-ranked unmasking & 40 & 39 & 0 & 1 \\ % src: analysis/census/gap-rulings.log (columns: ranks in F2 denom, non-plain, NS, U)
B-21 & combined letters; G or M with token rule unstated (+NS) & 19 & 18 & 1 & 0 \\
B-24 & sweep must be on the F1 task & 18 & 12 & 6 & 0 \\
B-10 & benchmark accuracy as generated-sample metric & 12 & 7 & 5 & 0 \\
B-9 & F1 from the in-scope task & 14 & 7 & 5 & 2 \\
B-19 & particle resampling as M & 6 & 6 & 0 & 0 \\
B-13 & conflicting statements about the rule & 3 & 3 & 0 & 0 \\
B-17 & best-of-$N$ and filtering as M & 3 & 3 & 0 & 0 \\
B-22 & diversity metrics not on the v1.1 list (decides F3) & 1 & 1 & 0 & 0 \\
\bottomrule
\end{tabular}
\end{table}

\setcounter{section}{4}
\section{The design-space table of the 47 models}
\label{sec:s-design}

\subsection{Condensed table and code key}
\label{sec:s-design-table}

Table~\ref{tab:design} is the condensed table of the 47 models (Section~\Mref{sec:design-table}); Figure~\Mref{fig:taxonomy} shows the same models as a dot matrix. The codes are as follows. \emph{C1}: Er erasure, Un uniform, Gs Gaussian on embeddings, Sx Gaussian on a logit simplex, Q general Markov or continuous-time transition matrices, Fl probability path (flow), Ed edit channel, AR+$X$ autoregressive text with diffusion component $X$ for images, +aux an auxiliary latent. \emph{Position rule (C4)}: rnd random sets, cnf ranked by confidence, ent entropy or threshold rule, blk blocks, all every position updated each step with no committed set, ar left to right, ns not stated. \emph{Value rule (C4)}: A argmax or greedy, T temperature, P top-$p$ or top-$k$, S sample-then-filter, G guidance, M sequence-level selection; plain sampled values; lik.\ the headline is a likelihood bound; ns not stated in the parts read; $\dagger$ temperature varied on purpose to trace an entropy frontier; (M) an optional variant. \emph{Revision (C5)}: none, impl implicit, corr corrector, remask, edit, out output-level, ns not stated. \emph{Levers}: L1 to L14 as in Table~\Mref{tab:levers}; SL sequence-level selection, for which the lever table has no row.
Position and revision codes in Table~\ref{tab:design} are condensed by us from the free-text cells of the long form.

{\scriptsize
\setlength{\tabcolsep}{3pt}
\begin{longtable}{@{}l p{3.0cm} l l l l p{1.9cm} l@{}}
\caption{Condensed design-space table of 47 models (M01 to M47). Codes as in the text. Rows M01 to M29 were copied from earlier readings; M30 to M47 were read at section level. ``partly'' marks the 11 stated limits of the notation (Section~\Mref{sec:design-limits}).}\label{tab:design}\\
\toprule
Id & Model & C1 & Position & Value & Revision & Levers & Fit \\
\midrule
\endfirsthead
\toprule
Id & Model & C1 & Position & Value & Revision & Levers & Fit \\
\midrule
\endhead
\bottomrule
\endfoot
M01 & D3PM~\cite{arXiv:2107.03006} & Q & rnd & lik. & impl & L1,L6 & yes \\
M02 & Campbell et al.~\cite{arXiv:2205.14987} & Q & rnd & ns & corr & L1,L9 & yes \\
M03 & SEDD~\cite{arXiv:2310.16834} & Er/Un & rnd & lik. & none/impl & L1,L6 & yes \\
M04 & MDLM~\cite{arXiv:2406.07524} & Er & rnd & lik. & none & L1,L4,L13 & yes \\
M05 & MD4~\cite{arXiv:2406.04329} & Er & rnd & lik. & none & L1,L6 & yes \\
M06 & RADD~\cite{arXiv:2406.03736} & Er & rnd & lik. & none & L6,L13 & yes \\
M07 & Zheng et al.~\cite{arXiv:2409.02908} & Er & rnd & T & none & L1,L11 & yes \\
M08 & ARDM~\cite{arXiv:2110.02037} & Er & rnd & lik. & none & L1,L6 & yes \\
M09 & LLaDA~\cite{arXiv:2502.09992} & Er & cnf & A & none & L2,L4,L11 & yes \\
M10 & MaskGIT~\cite{arXiv:2202.04200} & Er & cnf & S,T & none & L1,L11 & yes \\
M11 & ReMDM~\cite{arXiv:2503.00307} & Er & rnd & P & remask & L9,L11 & yes \\
M12 & EB-sampler~\cite{arXiv:2505.24857} & Er & ent & A & none & L1,L2 & yes \\
M13 & Diffusion-LM~\cite{arXiv:2205.14217} & Gs & all & A,G & impl & L1,L11,L12 & yes \\
M14 & CDCD~\cite{arXiv:2211.15089} & Gs & all & A & impl & L7,L11,L12 & yes \\
M15 & Plaid~\cite{arXiv:2305.18619} & Gs & all & A,T & impl & L7,L11,L12 & yes \\
M16 & DFM~\cite{arXiv:2407.15595} & Fl & rnd & ns & corr & L1,L9 & yes \\
M17 & Duo~\cite{arXiv:2506.10892} & Un & all & A & impl & L8,L9,L11 & yes \\
M18 & BD3-LM~\cite{arXiv:2503.09573} & Er & blk & P & none & L4,L13 & yes \\
M19 & SSD-LM~\cite{arXiv:2210.17432} & Sx & blk & A,P & impl & L4,L11 & yes \\
M20 & Show-o~\cite{arXiv:2408.12528} & AR+Er & ar, cnf & S,G & none & L5,L11,L12 & yes \\
M21 & UniDisc~\cite{arXiv:2503.20853} & Er & cnf & S,P,T,G & none & L5,L11,L12 & yes \\
M22 & MMaDA~\cite{arXiv:2505.15809} & Er & cnf & G & none & L2,L4,L12 & yes \\
M23 & MMaDA-Parallel~\cite{arXiv:2511.09611} & Er & cnf & G & none & L5,L12 & yes \\
M24 & FUDOKI~\cite{arXiv:2505.20147} & Fl & all & plain (M) & impl & L9 & \textbf{partly} \\
M25 & Lumina-DiMOO~\cite{arXiv:2510.06308} & Er & cnf & G & none & L2,L12,L14 & yes \\
M26 & Transfusion~\cite{arXiv:2408.11039} & AR+Gs & ar, all & G & impl & L5,L12 & \textbf{partly} \\
M27 & BAGEL~\cite{arXiv:2505.14683} & AR+Fl & ar, all & G & impl, out & L5,L10,L12 & \textbf{partly} \\
M28 & Show-o2~\cite{arXiv:2506.15564} & AR+Fl & ar, all & G & impl & L12 & \textbf{partly} \\
M29 & LaViDa-O~\cite{arXiv:2509.19244} & Er & ns & ns & out & L5,L10 & yes \\
\midrule
M30 & TESS~\cite{arXiv:2305.08379} & Sx & all & A & impl & L1,L7,L11 & yes \\
M31 & TESS 2~\cite{arXiv:2502.13917} & Sx & all & A,G & impl & L6,L7,L11,L12 & yes \\
M32 & DiffuSeq~\cite{arXiv:2210.08933} & Gs & all & A,M & impl & L1,L11,SL & \textbf{partly} \\
M33 & CANDI~\cite{arXiv:2510.22510} & Er+Gs & rnd & T$^\dagger$ & none & L1,L11,L12 & \textbf{partly} \\
M34 & Simplax~\cite{arXiv:2608.10615} & Un+aux & all & T$^\dagger$ & impl & L9,L11 & \textbf{partly} \\
M35 & Dream 7B~\cite{arXiv:2508.15487} & Er & ns & ns & ns & L2,L6 & yes \\
M36 & DiffuLLaMA~\cite{arXiv:2410.17891} & Er & rnd & P & none & L1,L6,L11 & yes \\
M37 & DiffuCoder~\cite{arXiv:2506.20639} & Er & cnf & T & none & L2,L6,L11 & yes \\
M38 & Seed Diffusion~\cite{arXiv:2508.02193} & Er+Ed & blk & ns & edit & L4,L6,L8,L14 & \textbf{partly} \\
M39 & Fast-dLLM~\cite{arXiv:2505.22618} & Er & cnf & A & none & L1,L2,L4,L14 & yes \\
M40 & Di4C~\cite{arXiv:2410.08709} & Er & rnd & plain & none & L3,L8 & \textbf{partly} \\
M41 & DCD~\cite{arXiv:2410.01949} & Er & blk & plain & none & L3 & \textbf{partly} \\
M42 & UDLM~\cite{arXiv:2412.10193} & Un & all & G & impl & L1,L9,L12 & yes \\
M43 & LaViDa~\cite{arXiv:2505.16839} & Er & ns & ns & none & L1,L6,L13 & yes \\
M44 & Muddit~\cite{arXiv:2505.23606} & Er & ns & G & none & L1,L12 & yes \\
M45 & Dynin-Omni~\cite{arXiv:2604.00007} & Er & cnf & S,G & none & L2,L4,L11,L12 & yes \\
M46 & D-DiT~\cite{arXiv:2501.00289} & Er+Fl & all, rnd & G & impl & L1,L12 & \textbf{partly} \\
M47 & Dimple~\cite{arXiv:2505.16990} & Er & cnf & A & none & L1,L2,L11 & yes \\

\end{longtable}
}

\subsection{Long form}
\label{sec:s-design-long}

\todo{long form of the design-space table with every cell and its source location (research/review/coverage/design-space-v1.csv); it is promised in Section~\Mref{sec:design-table}}

\setcounter{section}{6}
\section{Additional detail on the EXP-003 measurement}
\label{sec:s-exp003}

These four blocks are moved from Section~\Mref{sec:exp003}. The results they refer to are in Table~\Mref{tab:exp003-arms}.

\subsection{Deviations from the preregistered plan}
\label{sec:s-exp003-dev}

(1) An encoding error in the analysis and scoring code mis-keyed non-ASCII nouns; it was fixed before the controls were evaluated and before any main-stage sample existed, and the pilot analysis, which served for sizing only, may carry it. PC-1 then failed and the main stage continued as the card allows, under W2. (2) Eleven interpretive readings of the card and one erratum were accepted before scoring; the erratum corrects the mechanism paragraph, since the unconditional guidance branch masks every observed position and not only the prefix. (3) The hard stop on GPU time was raised from 7.0 to 8.5 hours, with the user's approval, during generation and before any main-stage sample was scored; the run used 8.02 hours against a plan of 6. (4) The pilot freeze was re-recorded for one analysis file, after generation and before scoring, for the encoding fix and a change of line endings. (5) A scoring-driver progress check that looked in the wrong directory stopped scoring falsely; it was fixed and scoring resumed. (6) The frozen analysis omitted five items, computed afterwards and labelled post-freeze: guardrail (iv), the composition contrast, the diversity intervals, the achieved detectable effect, and the counts of undersized placebo pools; three secondary contrasts used 5,000 bootstrap replicates, not 10,000. The gate audit verified the freeze hashes and the order of events, reproduced the confirmatory file byte for byte, and found no deviation that a main-stage outcome could have influenced. % src: decisions D0017, D0018, D0022, D0023, D0024, D0025; A sections 1.2 to 1.4 and 3; runs/EXP-003/clock.json

\subsection{Diversity of the guided arm (exploratory, post-freeze)}
\label{sec:s-exp003-div}

The hypothesis that the guided arm lowers diversity (E3-H4) was not tested by the frozen analysis. A post-freeze, exploratory computation (bias-centred percentile intervals) puts the text side in the predicted direction: unigram entropy of the completion words is lower for G than for P64 by 0.44 nats (95\% interval $-0.55$ to $-0.33$), and the Jaccard overlap of the two seeds' noun sets is higher by 0.054 (0.035 to 0.075). The image side runs the other way: the mean sibling-image cosine is 0.68 for G and 0.76 for P64, with no interval. We make no statement that guidance lowers diversity; the text-side direction is the one expected when guidance trades diversity for fidelity~\cite{arXiv:2207.12598}. % src: X item5_e3h4_diversity ("exploratory, post-freeze"); F arms.*.image_sibling_cosine; A section 6 item 8; N section 4 item 5

\subsection{Negative and positive controls, and guardrails}
\label{sec:s-exp003-ctrl}

NC-1 passed (link $-0.0003$, 95\% interval $-0.052$ to 0.050; 45 prefixes), NC-2 passed on G (0.002; $-0.016$ to 0.021) and on P64 (0.005; $-0.012$ to 0.023), and PC-2 passed (matched presence 0.77, other-caption 0.17; 113 nouns). Guardrails (i) to (iii) pass: the completion length of G is 0.79 times that of P64 (limits 0.5 to 2.0), no completion is empty, and G has 1,084 discordant nouns against 1,666 (ratio 0.65; limit 0.4). Guardrail (iv), near-uniform G images, was not computed by the frozen analysis; recomputed afterwards from the stored per-image flags it is 1 image of 276 (limit 0.15).
% src: runs/EXP-003/ctrl/controls_decision.json; F guardrails; X item1_guardrail_iv_near_uniform; A sections 4.1 and 6 item 9 (moved with the sentences from the "Checks and controls" paragraph of sections/09-exp003.tex, whose own src comment is unchanged)

\subsection{Statistic, inference, controls and checks (moved from Section~\Mref{sec:exp003})}
\label{sec:s-exp003-detail}
Moved unchanged from the main text for length.

\paragraph{Statistic.}
For prefix $p$ and seed $e\in\{0,1\}$ let $x_{pe}$ be the image and $D_{pe}$ the nouns that this completion names and the sibling completion (same prefix, other seed) does not, compared after removing a plural suffix. Let $y(n,x)\in\{0,1\}$ be the checker's answer to ``Is there a $n$ in this image?'', yes when $P(\text{yes})\ge 0.5$. A prefix is eligible when $D_{p0}$ and $D_{p1}$ are both non-empty. Within one arm its seed-balanced link is
\begin{equation}
  L_p \;=\; \tfrac12\Biggl[\frac{1}{|D_{p0}|}\sum_{n\in D_{p0}}\bigl(y(n,x_{p0})-y(n,x_{p1})\bigr)\;+\;\frac{1}{|D_{p1}|}\sum_{n\in D_{p1}}\bigl(y(n,x_{p1})-y(n,x_{p0})\bigr)\Biggr],
  \label{eq:s-exp003-link}
\end{equation}
and the link of arm $a$ is the mean over its eligible prefixes $\mathcal P_a$,
\begin{equation}
  \Lambda_{\mathrm b}(a) \;=\; \frac{1}{|\mathcal P_a|}\sum_{p\in\mathcal P_a} L_p^{(a)} .
  \label{eq:s-exp003-lambda}
\end{equation}
Equation~\eqref{eq:s-exp003-link} asks how much more often a noun that one seed's completion names is found in that seed's own image than in the sibling image. The sibling shares the prefix, the arm and the guidance setting, so what the prefix explains cancels; each seed has weight one half, so an additive yes-bias of either image cancels. The contrasts are $\Delta_1=\Lambda_{\mathrm b}(G)-\Lambda_{\mathrm b}(P64)$, $\Delta_{2a}=\Lambda_{\mathrm b}(G)-\Lambda_{\mathrm b}(O0)$ and $\Delta_{2b}=\Lambda_{\mathrm b}(O0)-\Lambda_{\mathrm b}(P32)$, each on the prefixes eligible in both its arms. On the prefixes eligible in all four arms,
\begin{equation}
  \Delta_1 \;=\; \Delta_{2a}+\Delta_{2b}+\bigl[\Lambda_{\mathrm b}(P32)-\Lambda_{\mathrm b}(P64)\bigr],
  \label{eq:s-exp003-identity}
\end{equation}
by telescoping; the bracket is the change from 32 to 64 plain passes. % src: C sections 4.1 and 4.2
The same-prefix sibling baseline is an instance of the same-condition independent-draw control used for co-generated outputs in MUNITE and CoDi~\cite{arXiv:2610.09866,arXiv:2305.11846}; seed balancing and the discordant-noun restriction are the parts we add. Yes/no presence probes follow POPE, TIFA and DSG~\cite{arXiv:2305.10355,arXiv:2303.11897,arXiv:2310.18235}, and judging joint coherence with a separate model is done as in MMVAE~\cite{arXiv:1911.03393}. % src: N section 4 items 4 and 5; decision D0019 (1)

\paragraph{Inference and controls.}
The prefix is the cluster; intervals are percentile bootstrap intervals with 10,000 replicates and the same resampled prefixes in every arm. The one primary test is $\Delta_1$ against zero, two-sided at 0.05; if it rejects, Holm's procedure at 0.05 covers $\Delta_{2a}$ and $\Delta_{2b}$. The minimum effect of interest, 0.10, is our choice and describes point estimates only. Four controls fix the wording: NC-1, a plain arm whose image ignores the completion given the prefix (link zero); NC-2, placebo nouns from another prefix; PC-1, a plain arm whose image is generated after the whole completion (link recoverable); PC-2, the checker on guided text-to-image images. A failed PC-1 turns every null label into ``not measurable with this instrument'' (rule W2) and leaves a positive readable. % src: C sections 4.5, 5, 6 and 7

\paragraph{Checks and controls.}
In the half of prefixes whose sibling images differ least in CLIP-IQA, $\Lambda_{\mathrm b}(G)$ is 0.131 (0.071 to 0.184; 63 prefixes), and the slope on the sibling quality difference is 0.160 ($-0.027$ to 0.366): the preregistered flag for quality-mediated coupling is not triggered, which does not exclude it. CLIP did not confirm the dependence: the sibling-link contrast of G against P64 on the completion text alone is 0.38 (100 times cosine similarity), 90\% interval $-0.085$ to 0.861 over 138 prefixes, where the rule requires a lower limit above zero. The composition sensitivity cannot be computed: G completions name more COCO-category nouns (share 0.110 against 0.065), but only 2 prefixes are eligible for it in both arms. NC-1, NC-2 and PC-2 passed (values in the supplement, Section~S7.3). PC-1 failed: the image-after-completion arm exceeds the independent arm by 0.031 (95\% interval $-0.042$ to 0.100; 44 prefixes), where 0.10 was required. We therefore cannot say that the checker registers a link even when the image sees the whole completion, and its sensitivity to within-prefix differences on guided images is untested. Guardrails (i) to (iii) pass, and guardrail (iv) was not computed by the frozen analysis; the values, with the post-freeze recomputation of (iv), are in the supplement (Section~S7.3). The plain arms' intervals in Table~\Mref{tab:exp003-arms} exclude zero. They are descriptive and uncorrected, and the statistic differs from the gate statistic of the earlier attempt on this model (Section~\Mref{sec:open-op3}). We therefore make no statement that plain sampling shows no link, and we do not read the intervals as replicating or contradicting that attempt. % src: F arms.G.quality_matched, clip_sibling_link_G_minus_P64_completion_only, W1_concordant, guardrails, arms.*.completion; runs/EXP-003/ctrl/controls_decision.json; X item1_guardrail_iv_near_uniform, item4_w3_composition; A sections 4.1 and 6 item 9; decision D0025 (5)


\setcounter{section}{7}
\section{Material moved from Sections 3, 4, 7 and 8}
\label{sec:s-moved}

Each block below is moved unchanged in content from the section named, where a one-sentence pointer remains.

\subsection{The commit-kernel sampler (from Section~\Mref{sec:notation-alg})}
\label{sec:s-alg}

Algorithm~\ref{alg:generate} is the sampler loop that calls Algorithm~\Mref{alg:commit} (Section~\Mref{sec:notation-alg}) once per step.

\begin{algorithm}[t]
\caption{\textsc{Generate}: the commit-kernel sampler}
\label{alg:generate}
\KwIn{predictor $p_\theta$; policy for \textsc{Commit} and a schedule $(k_n)_{n\le T}$; optional revision rule $\mathcal R$; conditioning $c$}
\KwOut{a sequence $\hat x\in\mathcal V^L$}
\KwData{state $s_n=(z_n,m_n)$ with memory $m_n$; the channel C1 fixes the initial state $z_0$}
$z_0\leftarrow$ all positions masked except those fixed by $c$; $m_0\leftarrow\varnothing$\;
\For{$n\leftarrow 1$ \KwTo $T$}{
  $(p_n,m_n)\leftarrow p_\theta(s_{n-1})$ \tcp*{one forward pass}
  $(A_n,y_n)\leftarrow\textsc{Commit}(s_{n-1},p_n,k_n)$ \tcp*{Algorithm~\Mref{alg:commit}; $k_n=|A_n|$}
  $z_n\leftarrow z_{n-1}$ with $y_n$ written at the positions $A_n$\;
  \If{$\mathcal R$ is given}{
    choose $R_n\subseteq\{1,\dots,L\}$ by $\mathcal R(z_n,p_n)$; set $z_n^{R_n}\leftarrow\mask$ \tcp*{C5, revision}
  }
}
\Return $z_T$\;
\end{algorithm}

\subsection{Rows marked ``yes'' that still changed the wording (from Section~\Mref{sec:design-limits})}
\label{sec:s-refine}

Four refinements were needed for rows that fit. Guidance by a classifier or reward gradient on logits or scores edits the state, not the marginals, so lever row L12 was renamed to cover it. The position rule often reads sampled or tempered values (MaskGIT, Show-o, UniDisc, Dynin-Omni, DiffuCoder), so position and value rules are separate rows in the table but coupled in practice. An exact cache can be bought by restricting the attention mask, which changes the learning term (LaViDa's prefix mask)~\cite{arXiv:2505.16839}. And the structured transition matrices already in D3PM need the same C1 entry as FUDOKI. % src: design-space-v1.md section 3(a) items 6 to 9

\subsection{Memory in the error budget (from Section~\Mref{sec:budget-other})}
\label{sec:s-memory}

Self-conditioning is a trained memory, not staleness, and belongs in $m_n$ and in $\Elearn$ on augmented states. With a memory the identity holds on the augmented path and is an equality for $x$ only if $m_n$ is a deterministic function of the revealed path; this is our inference, unchecked, and Simplax is a case where it fails (Section~\Mref{sec:design-limits}). % src: framework-v1.md section 2.1 (tagged [I]) and 2.5

\subsection{Withdrawn terms and statements (from Sections~\Mref{sec:budget-value} and~\Mref{sec:budget-props})}
\label{sec:s-withdrawn}

Terms (Section~\Mref{sec:budget-value}). An earlier version of our framework proposed two new terms, a selection bias for value-dependent rules and a tilting cost for temperature, truncation, argmax and guidance. Both are withdrawn (decision D0011). At one token per step the law of sample-then-filter equals that of a position-dependent state-only tilt in exact enumeration, and at more tokens per step it is a joint value law, so ``reads the sampled values'' describes an implementation and not a distribution. The claimed impossibility of per-step accounting of selection depended on one coupling, and the selection share is not canonical. What remains is published, and we credit it: the commit law at $k=1$ with heterogeneous competitors, its majorisation and the ``sample, select, resample'' fix~\cite{arXiv:2605.12522}; sampling then choosing against choosing then sampling, and the closeness of noised MaskGIT to a tilted sampler~\cite{arXiv:2510.04525}; confidence-based remasking against top-probability selection~\cite{arXiv:2609.34509}; the sample-then-filter and rank-then-sample vocabulary~\cite{arXiv:2604.00375}; and the one-token best-of-$N$ law with its $\log N-(N-1)/N$ bound~\cite{arXiv:2401.01879}. % src: decisions.jsonl D0011; framework-v2.md section A; dossier/prior-art-selection.md (comparison table; P1 to P5, EXTRA)

Statements (Section~\Mref{sec:budget-props}). That output entropy below a reference indicates value sharpening is false: with exact marginals and only factorisation error the entropy can fall far below the data entropy while the divergence stays small, so no reported entropy can by itself attribute a gain to a value rule, and attribution needs a controlled arm on the same checkpoint. The U-shaped dependence on the number of steps that a toy model produces is also not claimed (it appears in 20 of 32 toy variants and is sensitive to ties). % src: framework-v2.md section C (P2 withdrawn) and section D "Not claimed"; method/selection-term.md T4(b); red-team S4


\subsection{The cross-modal estimator (from Section~\Mref{sec:open-op3})}
\label{sec:s-estimator}

The estimator draws the text set from the model's joint ($k_{\mathrm T}$ sequential passes per sample, $R$ samples), re-queries the image marginals, and estimates $\hat I_j=\frac1R\sum_{r}\KL{p(b_j\mid s,a_r)}{\bar m_j}$, with $\bar m_j$ the mean of the $R$ re-queried marginals. The mixture is the reference because comparing with the one-pass marginal adds a marginal-inconsistency term that is non-zero even at zero dependence (0.081 nats per token in the red-team enumeration). The algebra was checked on toy models only and says nothing about re-query fidelity on a real network; the cost is $Rk_{\mathrm T}+R$ passes per probed step. % src: framework-v2.md section C (estimator paragraph, check_crossmodal_mixture.py); red-team block J

\bibliographystyle{ACM-Reference-Format}
\bibliography{references}

\end{document}
